<!doctype html><html lang="en"><head><meta charset="utf-8"><meta name="viewport" content="width=device-width,initial-scale=1"><title>Sealed source · paper_b_frozen.canonical.tex</title><style>
:root{color-scheme:light}*{box-sizing:border-box}html{scroll-behavior:smooth}body{background:#f8f7f3;color:#25342e;margin:0;font:14px/1.5 system-ui,sans-serif}.top{padding:23px max(22px,4vw);border-bottom:1px solid #d8dfd7;background:#fffefa}a{color:#2b6356}h1{font:normal clamp(24px,4vw,36px)/1.2 Georgia,serif;overflow-wrap:anywhere}.meta{font:12px/1.8 ui-monospace,monospace;color:#4b5d56;overflow-wrap:anywhere}.gate{border-left:3px solid #a66d3c;padding:12px 15px;max-width:850px;background:#f4ecdf;color:#594b39}.codewrap{overflow:auto;padding:18px max(16px,3vw)}pre{margin:0;min-width:max-content;font:12px/1.65 ui-monospace,SFMono-Regular,Consolas,monospace}.l{display:block;white-space:pre}.l:target{background:#f1e4c8;outline:2px solid #b79a6b}.n{display:inline-block;min-width:56px;padding-right:14px;text-align:right;color:#7e8b83;user-select:none;border-right:1px solid #e0e5dc;margin-right:15px}:focus-visible{outline:3px solid #ad7e35}@media(max-width:600px){.codewrap{padding:12px}pre{font-size:11px}}</style></head><body><header class="top"><p><a href="../../companion/technique_atlas/index.html">← Mathematical Machine</a> · <a href="../source_index.html">52 curated source waypoints</a></p><h1>paper_b_frozen.canonical.tex</h1><p class="meta">MANUSCRIPT/paper_b_frozen.canonical.tex · 1546 lines · SHA-256 9004c20d63218b529ea392d13914c17a347c485fd43814eb88bf8612a8512070</p><p class="gate"><b>Private evidence viewer.</b> This is byte-extracted from the sealed R0003 source archive. The source excerpts do not carry a newly approved public redistribution or software license. A hash and line match are not a fresh clean Lean replay.</p></header><div class="codewrap"><pre><span class="l" id="L1"><span class="n">1</span>\documentclass[11pt]{article}</span>
<span class="l" id="L2"><span class="n">2</span>\usepackage[margin=1.05in]{geometry}</span>
<span class="l" id="L3"><span class="n">3</span>\usepackage{amsmath,amssymb,amsthm,mathtools}</span>
<span class="l" id="L4"><span class="n">4</span>\usepackage{microtype}</span>
<span class="l" id="L5"><span class="n">5</span>\usepackage{booktabs}</span>
<span class="l" id="L6"><span class="n">6</span>\usepackage{enumitem}</span>
<span class="l" id="L7"><span class="n">7</span>\usepackage[hidelinks]{hyperref}</span>
<span class="l" id="L8"><span class="n">8</span>\usepackage[nameinlink,noabbrev]{cleveref}</span>
<span class="l" id="L9"><span class="n">9</span></span>
<span class="l" id="L10"><span class="n">10</span>\newtheorem{theorem}{Theorem}[section]</span>
<span class="l" id="L11"><span class="n">11</span>\newtheorem{proposition}[theorem]{Proposition}</span>
<span class="l" id="L12"><span class="n">12</span>\newtheorem{lemma}[theorem]{Lemma}</span>
<span class="l" id="L13"><span class="n">13</span>\newtheorem{corollary}[theorem]{Corollary}</span>
<span class="l" id="L14"><span class="n">14</span>\theoremstyle{definition}</span>
<span class="l" id="L15"><span class="n">15</span>\newtheorem{definition}[theorem]{Definition}</span>
<span class="l" id="L16"><span class="n">16</span>\theoremstyle{remark}</span>
<span class="l" id="L17"><span class="n">17</span>\newtheorem{remark}[theorem]{Remark}</span>
<span class="l" id="L18"><span class="n">18</span></span>
<span class="l" id="L19"><span class="n">19</span>\newcommand{\SSer}{\mathfrak S}</span>
<span class="l" id="L20"><span class="n">20</span>\newcommand{\Mob}{\mathsf M}</span>
<span class="l" id="L21"><span class="n">21</span>\newcommand{\ZE}{\mathsf Z_E}</span>
<span class="l" id="L22"><span class="n">22</span>\newcommand{\evE}{\operatorname{ev}_E}</span>
<span class="l" id="L23"><span class="n">23</span>\newcommand{\1}{\mathbf 1}</span>
<span class="l" id="L24"><span class="n">24</span>\newcommand{\e}{\mathrm e}</span>
<span class="l" id="L25"><span class="n">25</span>\newcommand{\ph}{\varphi}</span>
<span class="l" id="L26"><span class="n">26</span>\newcommand{\Om}{\Omega}</span>
<span class="l" id="L27"><span class="n">27</span>\newcommand{\Dcal}{\mathcal D}</span>
<span class="l" id="L28"><span class="n">28</span>\newcommand{\Fcal}{\mathcal F}</span>
<span class="l" id="L29"><span class="n">29</span>\newcommand{\Kcal}{\mathcal K}</span>
<span class="l" id="L30"><span class="n">30</span>\newcommand{\Lcal}{\mathcal L}</span>
<span class="l" id="L31"><span class="n">31</span>\newcommand{\Zmod}{\mathbb Z}</span>
<span class="l" id="L32"><span class="n">32</span>\newcommand{\eps}{\varepsilon}</span>
<span class="l" id="L33"><span class="n">33</span></span>
<span class="l" id="L34"><span class="n">34</span>\title{An Exact Contraction Dictionary and a Fixed-Modulus Anchored Transfer Theorem for Singular Series}</span>
<span class="l" id="L35"><span class="n">35</span>\author{Charlie Dorsey}</span>
<span class="l" id="L36"><span class="n">36</span>\date{September 2026}</span>
<span class="l" id="L37"><span class="n">37</span></span>
<span class="l" id="L38"><span class="n">38</span>\begin{document}</span>
<span class="l" id="L39"><span class="n">39</span>\maketitle</span>
<span class="l" id="L40"><span class="n">40</span></span>
<span class="l" id="L41"><span class="n">41</span>\begin{abstract}</span>
<span class="l" id="L42"><span class="n">42</span>Fix a squarefree modulus $q$ and compatible reduced endpoint classes $a,b\pmod q$.  We organize the endpoint-anchored and free connected singular-series coordinates by a common Boolean contraction calculus, identify the completed endpoint coordinate exactly with the free-label connected coordinate, and isolate the anchored-to-free defect as the top two-anchor connected contraction residual.  We then prove, pointwise along compatible heights,</span>
<span class="l" id="L43"><span class="n">43</span>\[</span>
<span class="l" id="L44"><span class="n">44</span>A_{1,q}(a,b;h)-B_{1,q}(a,b;h)=o_q(\log h).</span>
<span class="l" id="L45"><span class="n">45</span>\]</span>
<span class="l" id="L46"><span class="n">46</span>The fixed-$q$ proof isolates the finite-period changes in the pointwise $q=1$ argument: the low-prime and admissibility data form a fixed periodic decoration, the high-prime Boolean Fourier coefficients have unchanged absolute value, the exterior completion and reciprocal estimates incur only fixed-conductor constants, and the one-dimensional logarithmic main terms cancel before the nonprincipal periodic corrections.  We also record a finite-cutoff arithmetic diagonal-contraction identity for the active high-prime coefficient system.  No rank-uniform conclusion is asserted.</span>
<span class="l" id="L47"><span class="n">47</span>\end{abstract}</span>
<span class="l" id="L48"><span class="n">48</span></span>
<span class="l" id="L49"><span class="n">49</span>\section{Introduction and scope}\label{sec:intro}</span>
<span class="l" id="L50"><span class="n">50</span></span>
<span class="l" id="L51"><span class="n">51</span>For a finite set $H\subset\Zmod$ and a squarefree integer $q\ge1$, define the $q$-restricted singular series</span>
<span class="l" id="L52"><span class="n">52</span>\begin{equation}\label{eq:S_q}</span>
<span class="l" id="L53"><span class="n">53</span>\SSer_q(H)</span>
<span class="l" id="L54"><span class="n">54</span>=\prod_{p\nmid q}</span>
<span class="l" id="L55"><span class="n">55</span>\left(1-\frac{\nu_p(H)}p\right)</span>
<span class="l" id="L56"><span class="n">56</span>\left(1-\frac1p\right)^{-|H|},</span>
<span class="l" id="L57"><span class="n">57</span>\end{equation}</span>
<span class="l" id="L58"><span class="n">58</span>where $\nu_p(H)$ is the number of residue classes occupied by $H$ modulo $p$.  As usual, the product is zero when the tuple is inadmissible at an active prime.  Let</span>
<span class="l" id="L59"><span class="n">59</span>\begin{equation}\label{eq:S_q0}</span>
<span class="l" id="L60"><span class="n">60</span>\SSer_{q,0}(H)</span>
<span class="l" id="L61"><span class="n">61</span>=\sum_{J\subseteq H}(-1)^{|H|-|J|}\SSer_q(J)</span>
<span class="l" id="L62"><span class="n">62</span>\end{equation}</span>
<span class="l" id="L63"><span class="n">63</span>be its Boolean M\&quot;obius transform.</span>
<span class="l" id="L64"><span class="n">64</span></span>
<span class="l" id="L65"><span class="n">65</span>The purpose of this note is narrow.  First, we give an exact finite dictionary between endpoint completion and free-label connected coordinates.  Second, we prove the fixed-squarefree-$q$ pointwise level-one anchored-to-free transfer theorem.  Third, we record a finite-cutoff arithmetic diagonal-contraction identity for the active high-prime coefficient system.  We do not address an all-orders tail, a growing-rank theorem, or any resummation problem.</span>
<span class="l" id="L66"><span class="n">66</span></span>
<span class="l" id="L67"><span class="n">67</span>At level one, the two quantities compared by the theorem have simple roles.  The sum $A_{1,q}$ measures the connected contribution of one admissible interior point relative to the two fixed endpoints, while $B_{1,q}$ is the sum of the two one-anchor connected boundary contributions.  Their exact formulas appear in \eqref{eq:A1q}--\eqref{eq:B1q}; the theorem proves that their difference is $o_q(\log h)$ for every compatible height progression.</span>
<span class="l" id="L68"><span class="n">68</span></span>
<span class="l" id="L69"><span class="n">69</span>The proof has four stages.  The Boolean contraction identity in \cref{sec:dictionary} first identifies the correct connected coordinate exactly.  The fixed-period reduction in \cref{sec:fixed-period} then isolates the only changes caused by the modulus $q$: a bounded low-period decoration and the deletion of finitely many Euler factors.  The analytic work is concentrated in two places: the Boolean small box, where the companion $q=1$ argument imports Shiu&#x27;s short-interval estimate, and the exterior, where weighted completion, reciprocity, and Karatsuba&#x27;s trigonometric-sum estimate are used.  The remaining one-dimensional terms are matched by a fixed-conductor P\&#x27;etermann estimate.  Each imported result is stated below with the specific estimate and hypotheses used; the complete $q=1$ derivations are carried by the companion manuscript \cite{DorseyPointwise2026} rather than repeated here.</span>
<span class="l" id="L70"><span class="n">70</span></span>
<span class="l" id="L71"><span class="n">71</span>Fixed physical anchors and the conditional fixed-gap Poisson factor have established antecedents in the consecutive-prime literature \cite{GoldstonLedoan2011,LOS}; congruence-restricted singular-series sums likewise have substantial prior art \cite{Kuperberg2025}.  So do Boolean M\&quot;obius inversion and generic connected/cumulant language.  No novelty claim for those ingredients is made here.  The priority posture of the present theorem is stated separately in \cref{sec:priority}; mathematical statements below do not depend on that posture.</span>
<span class="l" id="L72"><span class="n">72</span></span>
<span class="l" id="L73"><span class="n">73</span>\section{The connected contraction dictionary}\label{sec:dictionary}</span>
<span class="l" id="L74"><span class="n">74</span></span>
<span class="l" id="L75"><span class="n">75</span>Fix the endpoint set</span>
<span class="l" id="L76"><span class="n">76</span>\[</span>
<span class="l" id="L77"><span class="n">77</span>E=\{0,h\}.</span>
<span class="l" id="L78"><span class="n">78</span>\]</span>
<span class="l" id="L79"><span class="n">79</span>Let $T$ be any finite set of free labels disjoint from $E$.  For $A\subseteq E$ and $U\subseteq T$, put</span>
<span class="l" id="L80"><span class="n">80</span>\[</span>
<span class="l" id="L81"><span class="n">81</span>f_q(A,U)=\SSer_q(A\cup U).</span>
<span class="l" id="L82"><span class="n">82</span>\]</span>
<span class="l" id="L83"><span class="n">83</span>Define partial Boolean M\&quot;obius operators</span>
<span class="l" id="L84"><span class="n">84</span>\begin{align}</span>
<span class="l" id="L85"><span class="n">85</span>(\Mob_E f)(A,U)</span>
<span class="l" id="L86"><span class="n">86</span>&amp;=\sum_{B\subseteq A}(-1)^{|A|-|B|}f(B,U),\label{eq:ME}\\</span>
<span class="l" id="L87"><span class="n">87</span>(\Mob_T f)(A,U)</span>
<span class="l" id="L88"><span class="n">88</span>&amp;=\sum_{V\subseteq U}(-1)^{|U|-|V|}f(A,V).\label{eq:MT}</span>
<span class="l" id="L89"><span class="n">89</span>\end{align}</span>
<span class="l" id="L90"><span class="n">90</span>They commute, and</span>
<span class="l" id="L91"><span class="n">91</span>\[</span>
<span class="l" id="L92"><span class="n">92</span>(\Mob_E\Mob_Tf_q)(A,U)=\SSer_{q,0}(A\cup U).</span>
<span class="l" id="L93"><span class="n">93</span>\]</span>
<span class="l" id="L94"><span class="n">94</span>Let endpoint-fiber contraction be</span>
<span class="l" id="L95"><span class="n">95</span>\[</span>
<span class="l" id="L96"><span class="n">96</span>(\ZE g)(U)=\sum_{A\subseteq E}g(A,U),</span>
<span class="l" id="L97"><span class="n">97</span>\]</span>
<span class="l" id="L98"><span class="n">98</span>and let $\evE$ denote evaluation at the full endpoint set $A=E$.</span>
<span class="l" id="L99"><span class="n">99</span></span>
<span class="l" id="L100"><span class="n">100</span>\begin{proposition}[Endpoint contraction identity]\label{prop:endpoint-contraction}</span>
<span class="l" id="L101"><span class="n">101</span>For every $q$, every endpoint pair $E$, and every finite free set $T$,</span>
<span class="l" id="L102"><span class="n">102</span>\begin{equation}\label{eq:commuting-identity}</span>
<span class="l" id="L103"><span class="n">103</span>\boxed{\ZE\Mob_E\Mob_T=\Mob_T\evE.}</span>
<span class="l" id="L104"><span class="n">104</span>\end{equation}</span>
<span class="l" id="L105"><span class="n">105</span>Equivalently,</span>
<span class="l" id="L106"><span class="n">106</span>\begin{equation}\label{eq:F_E}</span>
<span class="l" id="L107"><span class="n">107</span>\Fcal_{E,q}(T)</span>
<span class="l" id="L108"><span class="n">108</span>:=\sum_{A\subseteq E}\SSer_{q,0}(A\cup T)</span>
<span class="l" id="L109"><span class="n">109</span>=\Mob_T[\SSer_q(E\cup\cdot)](T).</span>
<span class="l" id="L110"><span class="n">110</span>\end{equation}</span>
<span class="l" id="L111"><span class="n">111</span>\end{proposition}</span>
<span class="l" id="L112"><span class="n">112</span></span>
<span class="l" id="L113"><span class="n">113</span>\begin{proof}</span>
<span class="l" id="L114"><span class="n">114</span>For fixed $U\subseteq T$,</span>
<span class="l" id="L115"><span class="n">115</span>\[</span>
<span class="l" id="L116"><span class="n">116</span>\sum_{A\subseteq E}\sum_{B\subseteq A}</span>
<span class="l" id="L117"><span class="n">117</span>(-1)^{|A|-|B|}f_q(B,U)</span>
<span class="l" id="L118"><span class="n">118</span>=f_q(E,U),</span>
<span class="l" id="L119"><span class="n">119</span>\]</span>
<span class="l" id="L120"><span class="n">120</span>which is the Boolean zeta--M\&quot;obius identity on the endpoint fiber.  Since $\Mob_E$ and $\Mob_T$ commute, \eqref{eq:commuting-identity} follows.</span>
<span class="l" id="L121"><span class="n">121</span>\end{proof}</span>
<span class="l" id="L122"><span class="n">122</span></span>
<span class="l" id="L123"><span class="n">123</span>Assume $\SSer_q(E)&gt;0$ and define</span>
<span class="l" id="L124"><span class="n">124</span>\begin{equation}\label{eq:Rq}</span>
<span class="l" id="L125"><span class="n">125</span>R_q(U)=\frac{\SSer_q(E\cup U)}{\SSer_q(E)},</span>
<span class="l" id="L126"><span class="n">126</span>\qquad</span>
<span class="l" id="L127"><span class="n">127</span>R_q^\circ(T)=\sum_{V\subseteq T}(-1)^{|T|-|V|}R_q(V).</span>
<span class="l" id="L128"><span class="n">128</span>\end{equation}</span>
<span class="l" id="L129"><span class="n">129</span>Then \cref{prop:endpoint-contraction} gives the pointwise identity</span>
<span class="l" id="L130"><span class="n">130</span>\begin{equation}\label{eq:F-Rcirc}</span>
<span class="l" id="L131"><span class="n">131</span>\boxed{\Fcal_{E,q}(T)=\SSer_q(E)R_q^\circ(T).}</span>
<span class="l" id="L132"><span class="n">132</span>\end{equation}</span>
<span class="l" id="L133"><span class="n">133</span>If $I$ is a finite interior set and</span>
<span class="l" id="L134"><span class="n">134</span>\[</span>
<span class="l" id="L135"><span class="n">135</span>C_{k,q}^{\mathrm{comp}}</span>
<span class="l" id="L136"><span class="n">136</span>=\sum_{\substack{T\subseteq I\\|T|=k}}\Fcal_{E,q}(T),</span>
<span class="l" id="L137"><span class="n">137</span>\qquad</span>
<span class="l" id="L138"><span class="n">138</span>A_{k,q}^{\circ}</span>
<span class="l" id="L139"><span class="n">139</span>=\sum_{\substack{T\subseteq I\\|T|=k}}R_q^\circ(T),</span>
<span class="l" id="L140"><span class="n">140</span>\]</span>
<span class="l" id="L141"><span class="n">141</span>then</span>
<span class="l" id="L142"><span class="n">142</span>\begin{equation}\label{eq:C-A-dictionary}</span>
<span class="l" id="L143"><span class="n">143</span>\boxed{C_{k,q}^{\mathrm{comp}}=\SSer_q(E)A_{k,q}^{\circ}.}</span>
<span class="l" id="L144"><span class="n">144</span>\end{equation}</span>
<span class="l" id="L145"><span class="n">145</span>At $q=1$ this is the exact normalization dictionary between the completed endpoint coordinate and the anchored free-label connected coordinate.</span>
<span class="l" id="L146"><span class="n">146</span></span>
<span class="l" id="L147"><span class="n">147</span>\subsection{Anchor subtraction}</span>
<span class="l" id="L148"><span class="n">148</span></span>
<span class="l" id="L149"><span class="n">149</span>Write</span>
<span class="l" id="L150"><span class="n">150</span>\[</span>
<span class="l" id="L151"><span class="n">151</span>c_q(A,T)=\SSer_{q,0}(A\cup T),</span>
<span class="l" id="L152"><span class="n">152</span>\]</span>
<span class="l" id="L153"><span class="n">153</span>and define the lower-anchor sum and the top-anchor sector by</span>
<span class="l" id="L154"><span class="n">154</span>\begin{align}</span>
<span class="l" id="L155"><span class="n">155</span>b_q(T)&amp;=c_q(\varnothing,T)+c_q(\{0\},T)+c_q(\{h\},T),\label{eq:bq}\\</span>
<span class="l" id="L156"><span class="n">156</span>g_q(T)&amp;=c_q(E,T)=\SSer_{q,0}(E\cup T).\label{eq:gq}</span>
<span class="l" id="L157"><span class="n">157</span>\end{align}</span>
<span class="l" id="L158"><span class="n">158</span>Then $\Fcal_{E,q}(T)=b_q(T)+g_q(T)$.  Since singleton singular series are one,</span>
<span class="l" id="L159"><span class="n">159</span>\[</span>
<span class="l" id="L160"><span class="n">160</span>\SSer_{q,0}(E)=\SSer_q(E)-1.</span>
<span class="l" id="L161"><span class="n">161</span>\]</span>
<span class="l" id="L162"><span class="n">162</span>Combining this with \eqref{eq:F-Rcirc} gives the exact residual identity</span>
<span class="l" id="L163"><span class="n">163</span>\begin{equation}\label{eq:anchor-residual}</span>
<span class="l" id="L164"><span class="n">164</span>\boxed{</span>
<span class="l" id="L165"><span class="n">165</span>R_q^\circ(T)-b_q(T)</span>
<span class="l" id="L166"><span class="n">166</span>=\frac{g_q(T)-\SSer_{q,0}(E)b_q(T)}{\SSer_q(E)}.</span>
<span class="l" id="L167"><span class="n">167</span>}</span>
<span class="l" id="L168"><span class="n">168</span>\end{equation}</span>
<span class="l" id="L169"><span class="n">169</span>Thus the anchored-to-free transfer defect is the normalized failure of the top two-anchor connected sector to contract onto the lower anchor sectors.</span>
<span class="l" id="L170"><span class="n">170</span></span>
<span class="l" id="L171"><span class="n">171</span>\subsection{The same transform at one prime}</span>
<span class="l" id="L172"><span class="n">172</span></span>
<span class="l" id="L173"><span class="n">173</span>At one active prime, write</span>
<span class="l" id="L174"><span class="n">174</span>\[</span>
<span class="l" id="L175"><span class="n">175</span>L_p(H)=\mathbf E_z\prod_{x\in H}\phi_z(x),</span>
<span class="l" id="L176"><span class="n">176</span>\]</span>
<span class="l" id="L177"><span class="n">177</span>where $\phi_z(x)$ is the normalized local avoidance factor.  If a coordinate value is repeated, $H$ is read as a labeled multiset and the product is over coordinate occurrences.  Boolean connectedization gives</span>
<span class="l" id="L178"><span class="n">178</span>\begin{equation}\label{eq:local-connected}</span>
<span class="l" id="L179"><span class="n">179</span>\sum_{J\subseteq H}(-1)^{|H|-|J|}L_p(J)</span>
<span class="l" id="L180"><span class="n">180</span>=\mathbf E_z\prod_{x\in H}(\phi_z(x)-1).</span>
<span class="l" id="L181"><span class="n">181</span>\end{equation}</span>
<span class="l" id="L182"><span class="n">182</span>This is the local-prime image of the same connected transform: every free Fourier frequency is forced away from zero after connectedization.</span>
<span class="l" id="L183"><span class="n">183</span></span>
<span class="l" id="L184"><span class="n">184</span>\section{Fixed squarefree modulus and the level-one theorem}\label{sec:fixed-q-setup}</span>
<span class="l" id="L185"><span class="n">185</span></span>
<span class="l" id="L186"><span class="n">186</span>Fix squarefree $q\ge1$ and endpoint classes</span>
<span class="l" id="L187"><span class="n">187</span>\[</span>
<span class="l" id="L188"><span class="n">188</span>a,b\in(\Zmod/q\Zmod)^\times,</span>
<span class="l" id="L189"><span class="n">189</span>\qquad h\equiv b-a\pmod q.</span>
<span class="l" id="L190"><span class="n">190</span>\]</span>
<span class="l" id="L191"><span class="n">191</span>Assume the endpoint pair is $q$-restricted admissible,</span>
<span class="l" id="L192"><span class="n">192</span>\[</span>
<span class="l" id="L193"><span class="n">193</span>\SSer_q(E)&gt;0,</span>
<span class="l" id="L194"><span class="n">194</span>\qquad E=\{0,h\}.</span>
<span class="l" id="L195"><span class="n">195</span>\]</span>
<span class="l" id="L196"><span class="n">196</span>Every legal height is even: if $2\nmid q$, admissibility at the active prime $2$ forces $h$ even; if $2\mid q$, the unit classes $a,b$ are both odd and hence their difference is even modulo $2$, so $h\equiv b-a\equiv0\pmod2$.</span>
<span class="l" id="L197"><span class="n">197</span></span>
<span class="l" id="L198"><span class="n">198</span>Define the admissible interior and its exact size by</span>
<span class="l" id="L199"><span class="n">199</span>\begin{equation}\label{eq:IqNq}</span>
<span class="l" id="L200"><span class="n">200</span>I_q(a;h)=\{0&lt;t&lt;h:(a+t,q)=1\},</span>
<span class="l" id="L201"><span class="n">201</span>\qquad N_q(a;h)=|I_q(a;h)|.</span>
<span class="l" id="L202"><span class="n">202</span>\end{equation}</span>
<span class="l" id="L203"><span class="n">203</span>The level-one connected and lower-anchor sums are</span>
<span class="l" id="L204"><span class="n">204</span>\begin{align}</span>
<span class="l" id="L205"><span class="n">205</span>A_{1,q}(a,b;h)</span>
<span class="l" id="L206"><span class="n">206</span>&amp;=\sum_{t\in I_q(a;h)}R_q^\circ(\{t\})\notag\\</span>
<span class="l" id="L207"><span class="n">207</span>&amp;=\sum_{t\in I_q(a;h)}</span>
<span class="l" id="L208"><span class="n">208</span>\left(</span>
<span class="l" id="L209"><span class="n">209</span>\frac{\SSer_q(\{0,t,h\})}{\SSer_q(\{0,h\})}-1</span>
<span class="l" id="L210"><span class="n">210</span>\right),\label{eq:A1q}\\</span>
<span class="l" id="L211"><span class="n">211</span>B_{1,q}(a,b;h)</span>
<span class="l" id="L212"><span class="n">212</span>&amp;=\sum_{t\in I_q(a;h)}b_q(\{t\})\notag\\</span>
<span class="l" id="L213"><span class="n">213</span>&amp;=\sum_{t\in I_q(a;h)}</span>
<span class="l" id="L214"><span class="n">214</span>\left[\SSer_{q,0}(\{t\})</span>
<span class="l" id="L215"><span class="n">215</span>+\SSer_{q,0}(\{0,t\})</span>
<span class="l" id="L216"><span class="n">216</span>+\SSer_{q,0}(\{h,t\})\right].\label{eq:B1q}</span>
<span class="l" id="L217"><span class="n">217</span>\end{align}</span>
<span class="l" id="L218"><span class="n">218</span>Since $\SSer_{q,0}(\{t\})=0$,</span>
<span class="l" id="L219"><span class="n">219</span>\begin{equation}\label{eq:B1q-simple}</span>
<span class="l" id="L220"><span class="n">220</span>B_{1,q}</span>
<span class="l" id="L221"><span class="n">221</span>=\sum_{t\in I_q(a;h)}(\SSer_q(\{0,t\})-1)</span>
<span class="l" id="L222"><span class="n">222</span>+\sum_{t\in I_q(a;h)}(\SSer_q(\{h,t\})-1).</span>
<span class="l" id="L223"><span class="n">223</span>\end{equation}</span>
<span class="l" id="L224"><span class="n">224</span>By \eqref{eq:anchor-residual},</span>
<span class="l" id="L225"><span class="n">225</span>\begin{equation}\label{eq:AminusB-connected}</span>
<span class="l" id="L226"><span class="n">226</span>A_{1,q}-B_{1,q}</span>
<span class="l" id="L227"><span class="n">227</span>=\frac1{\SSer_q(E)}</span>
<span class="l" id="L228"><span class="n">228</span>\sum_{t\in I_q(a;h)}</span>
<span class="l" id="L229"><span class="n">229</span>\left[g_q(\{t\})-\SSer_{q,0}(E)b_q(\{t\})\right].</span>
<span class="l" id="L230"><span class="n">230</span>\end{equation}</span>
<span class="l" id="L231"><span class="n">231</span></span>
<span class="l" id="L232"><span class="n">232</span>\begin{theorem}[Fixed-squarefree-$q$ pointwise level-one transfer]\label{thm:fixed-q}</span>
<span class="l" id="L233"><span class="n">233</span>For every fixed squarefree $q$, uniformly over the finitely many compatible admissible endpoint residue pairs $(a,b)$,</span>
<span class="l" id="L234"><span class="n">234</span>\begin{equation}\label{eq:main-theorem}</span>
<span class="l" id="L235"><span class="n">235</span>\boxed{A_{1,q}(a,b;h)-B_{1,q}(a,b;h)=o_q(\log h)}</span>
<span class="l" id="L236"><span class="n">236</span>\end{equation}</span>
<span class="l" id="L237"><span class="n">237</span>as $h\to\infty$ through positive compatible heights.</span>
<span class="l" id="L238"><span class="n">238</span>\end{theorem}</span>
<span class="l" id="L239"><span class="n">239</span></span>
<span class="l" id="L240"><span class="n">240</span>There is no density-one restriction in $h$.</span>
<span class="l" id="L241"><span class="n">241</span></span>
<span class="l" id="L242"><span class="n">242</span>\subsection{Exact density normalization}</span>
<span class="l" id="L243"><span class="n">243</span></span>
<span class="l" id="L244"><span class="n">244</span>The admissible count is exactly periodic.</span>
<span class="l" id="L245"><span class="n">245</span></span>
<span class="l" id="L246"><span class="n">246</span>\begin{lemma}[Exact periodic count]\label{lem:Nq}</span>
<span class="l" id="L247"><span class="n">247</span>For each compatible pair $(a,b)$,</span>
<span class="l" id="L248"><span class="n">248</span>\begin{equation}\label{eq:Nq-periodic}</span>
<span class="l" id="L249"><span class="n">249</span>\boxed{N_q(a;h)=\frac{\ph(q)}q h+\eps_q(a,b),}</span>
<span class="l" id="L250"><span class="n">250</span>\end{equation}</span>
<span class="l" id="L251"><span class="n">251</span>where $\eps_q(a,b)$ depends only on the fixed compatible residue data, not on the positive representative of the gap class.</span>
<span class="l" id="L252"><span class="n">252</span>\end{lemma}</span>
<span class="l" id="L253"><span class="n">253</span></span>
<span class="l" id="L254"><span class="n">254</span>\begin{proof}</span>
<span class="l" id="L255"><span class="n">255</span>If $h$ is increased by $q$ while preserving its residue class, the interval gains one complete period of $q$ translates and hence exactly $\ph(q)$ admissible interior positions.  Therefore $N_q(a;h)-\ph(q)h/q$ is constant on each compatible height progression.</span>
<span class="l" id="L256"><span class="n">256</span>\end{proof}</span>
<span class="l" id="L257"><span class="n">257</span></span>
<span class="l" id="L258"><span class="n">258</span>For the completed Lemke Oliver--Soundararajan normalization, the fixed-modulus density parameter is</span>
<span class="l" id="L259"><span class="n">259</span>\begin{equation}\label{eq:deltaq}</span>
<span class="l" id="L260"><span class="n">260</span>\delta_q(x)=\frac{q}{\ph(q)\log x}.</span>
<span class="l" id="L261"><span class="n">261</span>\end{equation}</span>
<span class="l" id="L262"><span class="n">262</span>Consequently</span>
<span class="l" id="L263"><span class="n">263</span>\[</span>
<span class="l" id="L264"><span class="n">264</span>\delta_q(x)\frac{\ph(q)}q=\frac1{\log x}.</span>
<span class="l" id="L265"><span class="n">265</span>\]</span>
<span class="l" id="L266"><span class="n">266</span>The $q/\ph(q)$ factor belongs to the ambient admissible-point density; it does not appear as an additional multiplier inside the transfer kernel \eqref{eq:main-theorem}.  We use $\Delta\mu_0$ only as notation for the completed level-one coefficient per admissible interior site,</span>
<span class="l" id="L267"><span class="n">267</span>\[</span>
<span class="l" id="L268"><span class="n">268</span>\Delta\mu_0:=\frac{C_{1,q}^{\mathrm{comp}}}{N_q(a;h)}.</span>
<span class="l" id="L269"><span class="n">269</span>\]</span>
<span class="l" id="L270"><span class="n">270</span>Thus \eqref{eq:F-Rcirc} gives at level one</span>
<span class="l" id="L271"><span class="n">271</span>\begin{equation}\label{eq:LOS-level-one}</span>
<span class="l" id="L272"><span class="n">272</span>C_{1,q}^{\mathrm{comp}}=\SSer_q(E)A_{1,q},</span>
<span class="l" id="L273"><span class="n">273</span>\qquad</span>
<span class="l" id="L274"><span class="n">274</span>\Delta\mu_0=\frac{\SSer_q(E)}{N_q(a;h)}A_{1,q}.</span>
<span class="l" id="L275"><span class="n">275</span>\end{equation}</span>
<span class="l" id="L276"><span class="n">276</span></span>
<span class="l" id="L277"><span class="n">277</span>\section{Exact fixed-period reduction}\label{sec:fixed-period}</span>
<span class="l" id="L278"><span class="n">278</span></span>
<span class="l" id="L279"><span class="n">279</span>The pointwise $q=1$ proof is carried by the companion estimate</span>
<span class="l" id="L280"><span class="n">280</span>\begin{equation}\label{eq:q1-input}</span>
<span class="l" id="L281"><span class="n">281</span>\sum_{0&lt;t&lt;h}\left(</span>
<span class="l" id="L282"><span class="n">282</span>\frac{\SSer(\{0,t,h\})}{\SSer(\{0,h\})}-1</span>
<span class="l" id="L283"><span class="n">283</span>\right)</span>
<span class="l" id="L284"><span class="n">284</span>-2\sum_{0&lt;t&lt;h}(\SSer(\{0,t\})-1)</span>
<span class="l" id="L285"><span class="n">285</span>=o(\log h)</span>
<span class="l" id="L286"><span class="n">286</span>\end{equation}</span>
<span class="l" id="L287"><span class="n">287</span>for even $h$.  We now isolate every change caused by fixed $q$ before invoking its analytic bounds.</span>
<span class="l" id="L288"><span class="n">288</span></span>
<span class="l" id="L289"><span class="n">289</span>Put</span>
<span class="l" id="L290"><span class="n">290</span>\[</span>
<span class="l" id="L291"><span class="n">291</span>\Lambda_q=\operatorname{lcm}(6,q),</span>
<span class="l" id="L292"><span class="n">292</span>\qquad</span>
<span class="l" id="L293"><span class="n">293</span>\operatorname{rad}_{\ge5}(n):=\prod_{\substack{p\mid n\\p\ge5}}p.</span>
<span class="l" id="L294"><span class="n">294</span>\]</span>
<span class="l" id="L295"><span class="n">295</span>We write $\e(t):=\exp(2\pi i t)$ and, for an integer modulus $M\ge1$,</span>
<span class="l" id="L296"><span class="n">296</span>$\e_M(x):=\e(x/M)$.  For the anchored ratio, collect all admissibility and active low-prime data into the bounded $\Lambda_q$-periodic weight</span>
<span class="l" id="L297"><span class="n">297</span>\begin{equation}\label{eq:lambdaq}</span>
<span class="l" id="L298"><span class="n">298</span>\lambda_{q,a,h}(y)</span>
<span class="l" id="L299"><span class="n">299</span>=\1_{(a+y,q)=1}</span>
<span class="l" id="L300"><span class="n">300</span>\prod_{\substack{p\in\{2,3\}\\p\nmid q}}</span>
<span class="l" id="L301"><span class="n">301</span>\frac{L_p(\{0,y,h\})}{L_p(\{0,h\})}.</span>
<span class="l" id="L302"><span class="n">302</span>\end{equation}</span>
<span class="l" id="L303"><span class="n">303</span>All remaining divisor variables are supported on primes $p\ge5$ with $p\nmid q$.</span>
<span class="l" id="L304"><span class="n">304</span></span>
<span class="l" id="L305"><span class="n">305</span>For $p\ge5$ set</span>
<span class="l" id="L306"><span class="n">306</span>\begin{equation}\label{eq:c0c1}</span>
<span class="l" id="L307"><span class="n">307</span>c_0(p)=\frac{p(p-3)}{(p-1)(p-2)},</span>
<span class="l" id="L308"><span class="n">308</span>\qquad</span>
<span class="l" id="L309"><span class="n">309</span>c_1(p)=\frac{p(p-2)}{(p-1)^2}.</span>
<span class="l" id="L310"><span class="n">310</span>\end{equation}</span>
<span class="l" id="L311"><span class="n">311</span>Define squarefree coefficient systems</span>
<span class="l" id="L312"><span class="n">312</span>\begin{align}</span>
<span class="l" id="L313"><span class="n">313</span>\omega_q(e)</span>
<span class="l" id="L314"><span class="n">314</span>&amp;=\mu^2(e)\1_{(e,6q)=1}\prod_{p\mid e}\frac1{p-3},\label{eq:omegaq}\\</span>
<span class="l" id="L315"><span class="n">315</span>\eta_q(e)</span>
<span class="l" id="L316"><span class="n">316</span>&amp;=\mu^2(e)\1_{(e,6q)=1}\prod_{p\mid e}\frac1{p-2},\label{eq:etaq}\\</span>
<span class="l" id="L317"><span class="n">317</span>v_{q,h}(e)</span>
<span class="l" id="L318"><span class="n">318</span>&amp;=\mu^2(e)\1_{(e,6q)=1}</span>
<span class="l" id="L319"><span class="n">319</span>\prod_{\substack{p\mid e\\p\nmid h}}\frac1{p-3}</span>
<span class="l" id="L320"><span class="n">320</span>\prod_{\substack{p\mid e\\p\mid h}}\frac1{p-2},\label{eq:vqh}\\</span>
<span class="l" id="L321"><span class="n">321</span>C_q^*(h)</span>
<span class="l" id="L322"><span class="n">322</span>&amp;=\prod_{\substack{p\ge5\\p\nmid qh}}c_0(p)</span>
<span class="l" id="L323"><span class="n">323</span>\prod_{\substack{p\ge5\\p\nmid q,\ p\mid h}}c_1(p).\label{eq:Cqstar}</span>
<span class="l" id="L324"><span class="n">324</span>\end{align}</span>
<span class="l" id="L325"><span class="n">325</span>For coprime high divisors, put</span>
<span class="l" id="L326"><span class="n">326</span>\begin{align}</span>
<span class="l" id="L327"><span class="n">327</span>N_{q,a}(h;e_1,e_2)</span>
<span class="l" id="L328"><span class="n">328</span>&amp;=\sum_{\substack{0&lt;y&lt;h\\e_1\mid y,\ e_2\mid h-y}}</span>
<span class="l" id="L329"><span class="n">329</span>\lambda_{q,a,h}(y),\label{eq:Nqea}\\</span>
<span class="l" id="L330"><span class="n">330</span>\Delta_{q,a}(h;e_1,e_2)</span>
<span class="l" id="L331"><span class="n">331</span>&amp;=N_{q,a}(h;e_1,e_2)-\frac{N_q(a;h)}{e_1e_2}.\label{eq:Deltaqea}</span>
<span class="l" id="L332"><span class="n">332</span>\end{align}</span>
<span class="l" id="L333"><span class="n">333</span></span>
<span class="l" id="L334"><span class="n">334</span>\begin{proposition}[Exact divisor expansion]\label{prop:divisor-expansion}</span>
<span class="l" id="L335"><span class="n">335</span>The anchored level-one sum admits the exact identity</span>
<span class="l" id="L336"><span class="n">336</span>\begin{equation}\label{eq:Tq-expansion}</span>
<span class="l" id="L337"><span class="n">337</span>\boxed{</span>
<span class="l" id="L338"><span class="n">338</span>A_{1,q}(a,b;h)</span>
<span class="l" id="L339"><span class="n">339</span>=C_q^*(h)</span>
<span class="l" id="L340"><span class="n">340</span>\sum_{\substack{e_1,e_2\ge1\\(e_1,e_2)=1,\ (e_2,h)=1}}</span>
<span class="l" id="L341"><span class="n">341</span>v_{q,h}(e_1)\omega_q(e_2)</span>
<span class="l" id="L342"><span class="n">342</span>\Delta_{q,a}(h;e_1,e_2).</span>
<span class="l" id="L343"><span class="n">343</span>}</span>
<span class="l" id="L344"><span class="n">344</span>\end{equation}</span>
<span class="l" id="L345"><span class="n">345</span>Every divisor on the right is supported on primes $p\ge5$ with $p\nmid q$.</span>
<span class="l" id="L346"><span class="n">346</span>\end{proposition}</span>
<span class="l" id="L347"><span class="n">347</span></span>
<span class="l" id="L348"><span class="n">348</span>\begin{proof}</span>
<span class="l" id="L349"><span class="n">349</span>Expand the active high-prime local ratio prime by prime.  Keep the fixed low-prime and admissibility factor \eqref{eq:lambdaq} unexpanded.  At each active high prime, the formal baseline together with its permitted selector or selectors sums to one.  Primes dividing $q$ occur in neither local singular-series factor nor the high divisor systems.  Thus insertion of the exact baseline $N_q(a;h)/(e_1e_2)$ reconstructs the literal subtraction $-N_q(a;h)$ in \eqref{eq:A1q}; no asymptotic density replacement is used.</span>
<span class="l" id="L350"><span class="n">350</span>\end{proof}</span>
<span class="l" id="L351"><span class="n">351</span></span>
<span class="l" id="L352"><span class="n">352</span>\subsection{Rowwise normalization}\label{sec:row-contract}</span>
<span class="l" id="L353"><span class="n">353</span>The two-dimensional estimates are applied row by row, and the actual row</span>
<span class="l" id="L354"><span class="n">354</span>height must remain visible.  We write $H$ for the parity-reduced kernel</span>
<span class="l" id="L355"><span class="n">355</span>height: $H=h/2$ on an ordinary row and $H=h/(2d)$ on a row carrying a</span>
<span class="l" id="L356"><span class="n">356</span>deterministic high divisor $d\mid h$.  Every progression root and oscillatory</span>
<span class="l" id="L357"><span class="n">357</span>phase in that row uses this $H$.  Let</span>
<span class="l" id="L358"><span class="n">358</span>$R\mid\operatorname{rad}_{\ge5}(h)$ be the squarefree product of additional</span>
<span class="l" id="L359"><span class="n">359</span>forbidden high height primes inherited from the row reduction, kept separate</span>
<span class="l" id="L360"><span class="n">360</span>from the fixed modulus $q$; active fixed-product moduli satisfy $(m,HqR)=1$.</span>
<span class="l" id="L361"><span class="n">361</span></span>
<span class="l" id="L362"><span class="n">362</span>For one exact low-period row, let $\lambda$ be its bounded periodic weight,</span>
<span class="l" id="L363"><span class="n">363</span>with period dividing $\Lambda_q$, and define</span>
<span class="l" id="L364"><span class="n">364</span>\begin{align}</span>
<span class="l" id="L365"><span class="n">365</span>N_\lambda(H)&amp;=\sum_{0&lt;y&lt;H}\lambda(y),\label{eq:Nlambda}\\</span>
<span class="l" id="L366"><span class="n">366</span>D_\lambda(e_1,e_2;H)</span>
<span class="l" id="L367"><span class="n">367</span>&amp;=\sum_{\substack{0&lt;y&lt;H\\e_1\mid y,\ e_2\mid H-y}}\lambda(y)</span>
<span class="l" id="L368"><span class="n">368</span> -\frac{N_\lambda(H)}{e_1e_2}.\label{eq:Dlambda}</span>
<span class="l" id="L369"><span class="n">369</span>\end{align}</span>
<span class="l" id="L370"><span class="n">370</span></span>
<span class="l" id="L371"><span class="n">371</span>\begin{lemma}[Exact fixed-$q$ row reduction]\label{lem:fixed-q-row-reduction}</span>
<span class="l" id="L372"><span class="n">372</span>For a row of \eqref{eq:Tq-expansion}, write</span>
<span class="l" id="L373"><span class="n">373</span>\[</span>
<span class="l" id="L374"><span class="n">374</span>d=(e_1,h),\qquad r=e_1/d.</span>
<span class="l" id="L375"><span class="n">375</span>\]</span>
<span class="l" id="L376"><span class="n">376</span>Thus $d$ is either $1$ or a squarefree divisor of</span>
<span class="l" id="L377"><span class="n">377</span>$\operatorname{rad}_{\ge5}(h)$ supported away from $q$, and put</span>
<span class="l" id="L378"><span class="n">378</span>\[</span>
<span class="l" id="L379"><span class="n">379</span>h_d=\frac hd=2H_d,</span>
<span class="l" id="L380"><span class="n">380</span>\qquad</span>
<span class="l" id="L381"><span class="n">381</span> a_d\equiv d^{-1}a\pmod q,</span>
<span class="l" id="L382"><span class="n">382</span>\]</span>
<span class="l" id="L383"><span class="n">383</span>and define the reduced low-period weight</span>
<span class="l" id="L384"><span class="n">384</span>\begin{equation}\label{eq:reduced-low-weight-q}</span>
<span class="l" id="L385"><span class="n">385</span>\lambda^{[d]}_{q,a,h}(z):=\lambda_{q,a_d,h_d}(2z).</span>
<span class="l" id="L386"><span class="n">386</span>\end{equation}</span>
<span class="l" id="L387"><span class="n">387</span>For every such row one has the exact identity</span>
<span class="l" id="L388"><span class="n">388</span>\begin{equation}\label{eq:fixed-q-row-reduction}</span>
<span class="l" id="L389"><span class="n">389</span>\boxed{</span>
<span class="l" id="L390"><span class="n">390</span>\Delta_{q,a}(h;dr,e_2)</span>
<span class="l" id="L391"><span class="n">391</span>=D_{\lambda^{[d]}_{q,a,h}}(r,e_2;H_d)</span>
<span class="l" id="L392"><span class="n">392</span>+\frac{\delta^{\rm per}_{q,a,h}(d)}{r e_2},</span>
<span class="l" id="L393"><span class="n">393</span>}</span>
<span class="l" id="L394"><span class="n">394</span>\end{equation}</span>
<span class="l" id="L395"><span class="n">395</span>where</span>
<span class="l" id="L396"><span class="n">396</span>\begin{equation}\label{eq:periodic-density-delta-q}</span>
<span class="l" id="L397"><span class="n">397</span>\delta^{\rm per}_{q,a,h}(d)</span>
<span class="l" id="L398"><span class="n">398</span>:=N_{\lambda^{[d]}_{q,a,h}}(H_d)-\frac{N_q(a;h)}d</span>
<span class="l" id="L399"><span class="n">399</span>=O_q(1)</span>
<span class="l" id="L400"><span class="n">400</span>\end{equation}</span>
<span class="l" id="L401"><span class="n">401</span>uniformly in the permitted $d$ and compatible heights.  Consequently the</span>
<span class="l" id="L402"><span class="n">402</span>periodic-density term in \eqref{eq:fixed-q-row-reduction} contributes</span>
<span class="l" id="L403"><span class="n">403</span>$O_q(1)$ on the ordinary two-dimensional family and</span>
<span class="l" id="L404"><span class="n">404</span>$O_q(\log\log h)$ after the outer decorated-row sum.</span>
<span class="l" id="L405"><span class="n">405</span>\end{lemma}</span>
<span class="l" id="L406"><span class="n">406</span></span>
<span class="l" id="L407"><span class="n">407</span>\begin{proof}</span>
<span class="l" id="L408"><span class="n">408</span>Write $y=dz$.  Since every prime of $d$ is at least $5$ and $d$ is a unit</span>
<span class="l" id="L409"><span class="n">409</span>modulo $6q$, multiplication by $d$ preserves the active $2,3$ local ratio and</span>
<span class="l" id="L410"><span class="n">410</span>sends the admissibility condition to $(a_d+z,q)=1$.  Hence, literally,</span>
<span class="l" id="L411"><span class="n">411</span>\[</span>
<span class="l" id="L412"><span class="n">412</span>\lambda_{q,a,h}(dz)=\lambda_{q,a_d,h_d}(z).</span>
<span class="l" id="L413"><span class="n">413</span>\]</span>
<span class="l" id="L414"><span class="n">414</span>The legal height $h_d$ is even.  If $2\nmid q$, its active $p=2$ factor kills</span>
<span class="l" id="L415"><span class="n">415</span>odd $z$; if $2\mid q$, the unit residue $a_d$ is odd and the admissibility</span>
<span class="l" id="L416"><span class="n">416</span>indicator does the same.  Thus $z=2x$ and, because $r,e_2$ are odd,</span>
<span class="l" id="L417"><span class="n">417</span>\[</span>
<span class="l" id="L418"><span class="n">418</span>\sum_{\substack{0&lt;y&lt;h\\dr\mid y,\ e_2\mid h-y}}</span>
<span class="l" id="L419"><span class="n">419</span>\lambda_{q,a,h}(y)</span>
<span class="l" id="L420"><span class="n">420</span>=</span>
<span class="l" id="L421"><span class="n">421</span>\sum_{\substack{0&lt;x&lt;H_d\\r\mid x,\ e_2\mid H_d-x}}</span>
<span class="l" id="L422"><span class="n">422</span>\lambda^{[d]}_{q,a,h}(x).</span>
<span class="l" id="L423"><span class="n">423</span>\]</span>
<span class="l" id="L424"><span class="n">424</span>Subtracting $N_q(a;h)/(dr e_2)$ on the left and</span>
<span class="l" id="L425"><span class="n">425</span>$N_{\lambda^{[d]}_{q,a,h}}(H_d)/(r e_2)$ on the right proves</span>
<span class="l" id="L426"><span class="n">426</span>\eqref{eq:fixed-q-row-reduction}.</span>
<span class="l" id="L427"><span class="n">427</span></span>
<span class="l" id="L428"><span class="n">428</span>For the bound, one complete $\Lambda_q$-period of</span>
<span class="l" id="L429"><span class="n">429</span>$z\mapsto\lambda_{q,a_d,h_d}(2z)$ has mean $2\ph(q)/q$: the admissibility</span>
<span class="l" id="L430"><span class="n">430</span>indicator contributes $\ph(q)/q$, and each active low-prime local ratio has</span>
<span class="l" id="L431"><span class="n">431</span>complete-period mean one before the parity compression.  Hence</span>
<span class="l" id="L432"><span class="n">432</span>\[</span>
<span class="l" id="L433"><span class="n">433</span>N_{\lambda^{[d]}_{q,a,h}}(H_d)</span>
<span class="l" id="L434"><span class="n">434</span>=\frac{\ph(q)}q h_d+O_q(1).</span>
<span class="l" id="L435"><span class="n">435</span>\]</span>
<span class="l" id="L436"><span class="n">436</span>By \eqref{eq:Nq-periodic},</span>
<span class="l" id="L437"><span class="n">437</span>$N_q(a;h)/d=(\ph(q)/q)h_d+O_q(1)$, proving</span>
<span class="l" id="L438"><span class="n">438</span>\eqref{eq:periodic-density-delta-q}.  For a fixed product $m=r e_2$, the</span>
<span class="l" id="L439"><span class="n">439</span>correction is independent of its Boolean partition and therefore costs at</span>
<span class="l" id="L440"><span class="n">440</span>most a constant multiple of</span>
<span class="l" id="L441"><span class="n">441</span>\[</span>
<span class="l" id="L442"><span class="n">442</span>\sum_m \omega_q(m)\frac{2^{\Om(m)}}m&lt;\infty.</span>
<span class="l" id="L443"><span class="n">443</span>\]</span>
<span class="l" id="L444"><span class="n">444</span>On decorated rows the remaining outer factor is bounded by</span>
<span class="l" id="L445"><span class="n">445</span>\[</span>
<span class="l" id="L446"><span class="n">446</span>\sum_{d\mid \operatorname{rad}_{\ge5}(h),\ (d,q)=1}v_{q,h}(d)</span>
<span class="l" id="L447"><span class="n">447</span>=\prod_{\substack{p\mid h\\p\ge5,\ p\nmid q}}</span>
<span class="l" id="L448"><span class="n">448</span>\left(1+\frac1{p-2}\right)</span>
<span class="l" id="L449"><span class="n">449</span>\ll_q\log\log h.</span>
<span class="l" id="L450"><span class="n">450</span>\]</span>
<span class="l" id="L451"><span class="n">451</span>This gives the stated correction ledger.</span>
<span class="l" id="L452"><span class="n">452</span>\end{proof}</span>
<span class="l" id="L453"><span class="n">453</span></span>
<span class="l" id="L454"><span class="n">454</span>For squarefree $m$ put</span>
<span class="l" id="L455"><span class="n">455</span>\begin{equation}\label{eq:Dcal-row}</span>
<span class="l" id="L456"><span class="n">456</span>\boxed{</span>
<span class="l" id="L457"><span class="n">457</span>\Dcal_{m,q}^{(\lambda)}(H)</span>
<span class="l" id="L458"><span class="n">458</span>=2^{-\Om(m)}\sum_{e_1e_2=m}D_\lambda(e_1,e_2;H).</span>
<span class="l" id="L459"><span class="n">459</span>}</span>
<span class="l" id="L460"><span class="n">460</span>\end{equation}</span>
<span class="l" id="L461"><span class="n">461</span>There are only $O_q(1)$ low rows, and all estimates below are uniform over</span>
<span class="l" id="L462"><span class="n">462</span>that finite family, so the superscript $(\lambda)$ may be suppressed after</span>
<span class="l" id="L463"><span class="n">463</span>this declaration.  Deleting the improper endpoint partitions contributes a</span>
<span class="l" id="L464"><span class="n">464</span>separate endpoint correction and is not part of the definition</span>
<span class="l" id="L465"><span class="n">465</span>\eqref{eq:Dcal-row}.  The high-prime weight used with this object is always</span>
<span class="l" id="L466"><span class="n">466</span>$\omega_q$ from \eqref{eq:omegaq}.</span>
<span class="l" id="L467"><span class="n">467</span></span>
<span class="l" id="L468"><span class="n">468</span>\subsection{Boolean fixed-product fibers}</span>
<span class="l" id="L469"><span class="n">469</span></span>
<span class="l" id="L470"><span class="n">470</span>Let $m=e_1e_2$ be squarefree, coprime to $\Lambda_q$, and let $u\bmod m$ be the CRT idempotent selecting the ordered partition $(e_1,e_2)$.  For a fixed low residue $c\bmod\Lambda_q$, the combined progression root is the unique class $t_{u,c}\bmod \Lambda_qm$ satisfying</span>
<span class="l" id="L471"><span class="n">471</span>\begin{equation}\label{eq:crt-root}</span>
<span class="l" id="L472"><span class="n">472</span>t_{u,c}\equiv Hu\pmod m,</span>
<span class="l" id="L473"><span class="n">473</span>\qquad</span>
<span class="l" id="L474"><span class="n">474</span>t_{u,c}\equiv c\pmod{\Lambda_q},</span>
<span class="l" id="L475"><span class="n">475</span>\end{equation}</span>
<span class="l" id="L476"><span class="n">476</span>where $H$ denotes the current reduced kernel height.  Write $\bar\Lambda$ for the inverse of $\Lambda_q\pmod m$ and $\bar m$ for the inverse of $m\pmod{\Lambda_q}$, and define</span>
<span class="l" id="L477"><span class="n">477</span>\[</span>
<span class="l" id="L478"><span class="n">478</span>W(t;m)=\frac1{2^{\Om(m)}}</span>
<span class="l" id="L479"><span class="n">479</span>\sum_{u^2\equiv u\ (m)}\e\!\left(\frac{tu}{m}\right).</span>
<span class="l" id="L480"><span class="n">480</span>\]</span>
<span class="l" id="L481"><span class="n">481</span>CRT gives the exact factorization</span>
<span class="l" id="L482"><span class="n">482</span>\begin{equation}\label{eq:fixed-q-fourier-factorization}</span>
<span class="l" id="L483"><span class="n">483</span>\boxed{</span>
<span class="l" id="L484"><span class="n">484</span>\frac1{2^{\Om(m)}}</span>
<span class="l" id="L485"><span class="n">485</span>\sum_{u^2\equiv u\ (m)}</span>
<span class="l" id="L486"><span class="n">486</span>\e\!\left(\frac{jt_{u,c}}{\Lambda_qm}\right)</span>
<span class="l" id="L487"><span class="n">487</span>=\e\!\left(\frac{jc\bar m}{\Lambda_q}\right)</span>
<span class="l" id="L488"><span class="n">488</span>W(jH\bar\Lambda;m).</span>
<span class="l" id="L489"><span class="n">489</span>}</span>
<span class="l" id="L490"><span class="n">490</span>\end{equation}</span>
<span class="l" id="L491"><span class="n">491</span>Hence</span>
<span class="l" id="L492"><span class="n">492</span>\begin{equation}\label{eq:absolute-fourier}</span>
<span class="l" id="L493"><span class="n">493</span>\left|</span>
<span class="l" id="L494"><span class="n">494</span>\frac1{2^{\Om(m)}}</span>
<span class="l" id="L495"><span class="n">495</span>\sum_u\e\!\left(\frac{jt_{u,c}}{\Lambda_qm}\right)</span>
<span class="l" id="L496"><span class="n">496</span>\right|</span>
<span class="l" id="L497"><span class="n">497</span>=|W(jH\bar\Lambda;m)|.</span>
<span class="l" id="L498"><span class="n">498</span>\end{equation}</span>
<span class="l" id="L499"><span class="n">499</span>The low residue contributes only a unit phase and multiplication by $\bar\Lambda$ permutes the high unit classes.</span>
<span class="l" id="L500"><span class="n">500</span></span>
<span class="l" id="L501"><span class="n">501</span>Passing from idempotents to sign roots $x=2u-1$, define</span>
<span class="l" id="L502"><span class="n">502</span>\[</span>
<span class="l" id="L503"><span class="n">503</span>V(t;m)=\frac1{2^{\Om(m)}}</span>
<span class="l" id="L504"><span class="n">504</span>\sum_{x^2\equiv1\ (m)}\e\!\left(\frac{tx}{m}\right).</span>
<span class="l" id="L505"><span class="n">505</span>\]</span>
<span class="l" id="L506"><span class="n">506</span>The complete-unit second moment from the $q=1$ proof therefore survives unchanged:</span>
<span class="l" id="L507"><span class="n">507</span>\begin{equation}\label{eq:second-moment}</span>
<span class="l" id="L508"><span class="n">508</span>\frac1{\ph(s)}</span>
<span class="l" id="L509"><span class="n">509</span>\sum_{x\in(\Zmod/s\Zmod)^\times}|V(Hjx;s)|^2</span>
<span class="l" id="L510"><span class="n">510</span>=\prod_{\substack{p\mid s\\p\nmid j}}</span>
<span class="l" id="L511"><span class="n">511</span>\frac{p-2}{2(p-1)}.</span>
<span class="l" id="L512"><span class="n">512</span>\end{equation}</span>
<span class="l" id="L513"><span class="n">513</span>This is the key fixed-period observation: fixed $q$ changes the phases, but not the absolute Boolean Fourier coefficients seen by the small box.</span>
<span class="l" id="L514"><span class="n">514</span></span>
<span class="l" id="L515"><span class="n">515</span>\section{Fixed-modulus analytic estimates}\label{sec:analytic-port}</span>
<span class="l" id="L516"><span class="n">516</span></span>
<span class="l" id="L517"><span class="n">517</span>The exact divisor expansion has now reduced the theorem to three analytic estimates.  The small box must retain the Boolean gain that makes its logarithmic exponent strictly less than one; the exterior must preserve the weighted-completion and reciprocal savings after the fixed conductor is adjoined; and the one-dimensional rows must have matching logarithmic main terms.  The complete $q=1$ arguments are proved in the companion manuscript \cite{DorseyPointwise2026}.  Here we state the precise fixed-$q$ changes needed to reuse those arguments.  Throughout, $q$ is fixed, so implicit constants may depend on it.</span>
<span class="l" id="L518"><span class="n">518</span></span>
<span class="l" id="L519"><span class="n">519</span>\subsection{The Boolean small box}</span>
<span class="l" id="L520"><span class="n">520</span></span>
<span class="l" id="L521"><span class="n">521</span>The $q=1$ proof splits a squarefree fixed product $m=sr$ into a smooth part and a rough part, applies Shiu&#x27;s short-interval Brun--Titchmarsh theorem for nonnegative multiplicative functions \cite[Chapter~3, Theorem~1]{ShiuThesis1980} to the rough factor in a reduced class modulo $s$, and combines \eqref{eq:second-moment} with weighted Cauchy.  In the range used here, Shiu supplies an upper bound for a nonnegative multiplicative function on a dyadic interval in one reduced class modulo $s$, with $s&lt;y^{3/4}$ and interval length $y&gt;x^{1/4}$; the moving rough functions have uniformly bounded prime values and vanish on higher prime powers.  Those are the only features of the theorem used below.</span>
<span class="l" id="L522"><span class="n">522</span></span>
<span class="l" id="L523"><span class="n">523</span>The fixed period $\Lambda_q$ does not enlarge the moving Shiu modulus.  To see this explicitly, expand the bounded periodic weight $\lambda$ by the finitely many residue classes $c\pmod{\Lambda_q}$.  For a dyadic product block $m\asymp M$, write uniquely $m=sr$, with every prime of $s$ at most</span>
<span class="l" id="L524"><span class="n">524</span>\[</span>
<span class="l" id="L525"><span class="n">525</span>Y=\exp((\log M)^{0.99})</span>
<span class="l" id="L526"><span class="n">526</span>\]</span>
<span class="l" id="L527"><span class="n">527</span>and every prime of $r$ larger than $Y$.  On the good range $s\le M^{0.01}$, \eqref{eq:fixed-q-fourier-factorization} and multiplicativity of the Boolean transform give</span>
<span class="l" id="L528"><span class="n">528</span>\begin{equation}\label{eq:fixed-q-shiu-factor}</span>
<span class="l" id="L529"><span class="n">529</span>|W(jH\bar\Lambda;sr)|^2</span>
<span class="l" id="L530"><span class="n">530</span>=|W(jH\bar\Lambda\bar r;s)|^2</span>
<span class="l" id="L531"><span class="n">531</span> |W(jH\bar\Lambda\bar s;r)|^2</span>
<span class="l" id="L532"><span class="n">532</span>\le |W(jH\bar\Lambda\bar r;s)|^2.</span>
<span class="l" id="L533"><span class="n">533</span>\end{equation}</span>
<span class="l" id="L534"><span class="n">534</span>Grouping the rough factor by its reduced residue class $r\pmod s$ therefore produces exactly the same progression problem modulo $s$ as in the companion Shiu calculation: multiplication by the fixed unit $\bar\Lambda$ merely permutes those classes.  The residue $c\pmod{\Lambda_q}$ contributes only the unit phase in \eqref{eq:fixed-q-fourier-factorization}, and there are $O_q(1)$ such residues.  Thus the rough-factor estimate, the complete-unit second moment \eqref{eq:second-moment}, and the weighted-Cauchy exponent are unchanged.  Deleting the finitely many primes dividing $q$ only removes nonnegative terms from the Shiu majorants.</span>
<span class="l" id="L535"><span class="n">535</span></span>
<span class="l" id="L536"><span class="n">536</span>The only additional counting error comes from replacing the exact weighted root count in \eqref{eq:Dlambda} by its complete-period mean.  For each fixed product and low residue this endpoint error is $O_q(1)$; after the Boolean weights are summed it contributes</span>
<span class="l" id="L537"><span class="n">537</span>\[</span>
<span class="l" id="L538"><span class="n">538</span>O_q(1)+O_q(\log Z),</span>
<span class="l" id="L539"><span class="n">539</span>\]</span>
<span class="l" id="L540"><span class="n">540</span>as recorded in \eqref{eq:periodic-defect-cost}.  This is $o(\log h)$ at the chosen scale.  We may therefore import the companion small-box exponent without any hidden moving-modulus cost.</span>
<span class="l" id="L541"><span class="n">541</span></span>
<span class="l" id="L542"><span class="n">542</span>With the same parameter choices</span>
<span class="l" id="L543"><span class="n">543</span>\[</span>
<span class="l" id="L544"><span class="n">544</span>\sigma=\frac23+0.01,</span>
<span class="l" id="L545"><span class="n">545</span>\qquad</span>
<span class="l" id="L546"><span class="n">546</span>\log Z=4(\log h)^\sigma(1+o(1)),</span>
<span class="l" id="L547"><span class="n">547</span>\]</span>
<span class="l" id="L548"><span class="n">548</span>the fixed-$q$ small-box bound is, for a fixed exponent $K_q$,</span>
<span class="l" id="L549"><span class="n">549</span>\begin{equation}\label{eq:fixed-q-smallbox}</span>
<span class="l" id="L550"><span class="n">550</span>\sum_{\substack{1&lt;m\le Z\\(m,HqR)=1}}</span>
<span class="l" id="L551"><span class="n">551</span>\omega_q(m)2^{\Om(m)}|\Dcal_{m,q}(H)|</span>
<span class="l" id="L552"><span class="n">552</span>\ll_q</span>
<span class="l" id="L553"><span class="n">553</span>(\log\log Z)^{K_q}(\log Z)^{2\sqrt{0.505}}</span>
<span class="l" id="L554"><span class="n">554</span>+O_q(\log Z),</span>
<span class="l" id="L555"><span class="n">555</span>\end{equation}</span>
<span class="l" id="L556"><span class="n">556</span>where $\Dcal_{m,q}$ is the corresponding centered Boolean discrepancy.  Since</span>
<span class="l" id="L557"><span class="n">557</span>\begin{equation}\label{eq:exponent}</span>
<span class="l" id="L558"><span class="n">558</span>\sigma\,2\sqrt{0.505}</span>
<span class="l" id="L559"><span class="n">559</span>=0.9617240306403449\ldots&lt;1,</span>
<span class="l" id="L560"><span class="n">560</span>\end{equation}</span>
<span class="l" id="L561"><span class="n">561</span>the centered fixed-$q$ small box is $o_q(\log h)$ pointwise.</span>
<span class="l" id="L562"><span class="n">562</span></span>
<span class="l" id="L563"><span class="n">563</span>The exact use of $N_q(a;h)$ instead of its mean is the correction isolated in</span>
<span class="l" id="L564"><span class="n">564</span>\cref{lem:fixed-q-row-reduction}; together with the endpoint deletion it costs only</span>
<span class="l" id="L565"><span class="n">565</span>\begin{equation}\label{eq:periodic-defect-cost}</span>
<span class="l" id="L566"><span class="n">566</span>O_q(1)+O_q(\log Z)=o(\log h),</span>
<span class="l" id="L567"><span class="n">567</span>\end{equation}</span>
<span class="l" id="L568"><span class="n">568</span>because the fixed-product endpoint defect is $O_q(1)$ and</span>
<span class="l" id="L569"><span class="n">569</span>\[</span>
<span class="l" id="L570"><span class="n">570</span>\sum_m\omega_q(m)\frac{2^{\Om(m)}}m&lt;\infty.</span>
<span class="l" id="L571"><span class="n">571</span>\]</span>
<span class="l" id="L572"><span class="n">572</span></span>
<span class="l" id="L573"><span class="n">573</span>\subsection{Exterior completion and reciprocity}</span>
<span class="l" id="L574"><span class="n">574</span></span>
<span class="l" id="L575"><span class="n">575</span>For each low class $c\bmod\Lambda_q$, the extra factor</span>
<span class="l" id="L576"><span class="n">576</span>\[</span>
<span class="l" id="L577"><span class="n">577</span>\e_{\Lambda_q}(jc\bar m)</span>
<span class="l" id="L578"><span class="n">578</span>\]</span>
<span class="l" id="L579"><span class="n">579</span>is a fixed-conductor periodic function of the high modulus.  On the unit classes modulo $\Lambda_q$ it decomposes into finitely many fixed character/additive-frequency components.</span>
<span class="l" id="L580"><span class="n">580</span></span>
<span class="l" id="L581"><span class="n">581</span>The weighted-completion input uses a unit-preserving squarefree-removal split.</span>
<span class="l" id="L582"><span class="n">582</span>The divisor $d$ comes from the companion coefficient $\gamma(d)$, so $d$ is</span>
<span class="l" id="L583"><span class="n">583</span>squarefree and $(d,6)=1$; the active support further gives $(d,qR)=1$.  Under</span>
<span class="l" id="L584"><span class="n">584</span>precisely these conditions,</span>
<span class="l" id="L585"><span class="n">585</span>\begin{equation}\label{eq:fixed-q-squarefree-removal}</span>
<span class="l" id="L586"><span class="n">586</span>\boxed{</span>
<span class="l" id="L587"><span class="n">587</span> \mu^2(dn)\1_{(dn,6R)=1}\1_{(dn,q)=1}</span>
<span class="l" id="L588"><span class="n">588</span> =\1_{(n,6q)=1}\mu^2(n)\1_{(n,dR)=1}</span>
<span class="l" id="L589"><span class="n">589</span>}\qquad</span>
<span class="l" id="L590"><span class="n">590</span>\bigl[d\ \text{squarefree},\ (d,6qR)=1\bigr].</span>
<span class="l" id="L591"><span class="n">591</span>\end{equation}</span>
<span class="l" id="L592"><span class="n">592</span>Only the high-prime factor $dR$ is expanded by inclusion--exclusion:</span>
<span class="l" id="L593"><span class="n">593</span>\[</span>
<span class="l" id="L594"><span class="n">594</span> \1_{(n,6q)=1}\mu^2(n)\1_{(n,dR)=1}</span>
<span class="l" id="L595"><span class="n">595</span> =\1_{(n,6q)=1}</span>
<span class="l" id="L596"><span class="n">596</span> \left(\sum_{\substack{c\mid dR\\c\mid n}}\mu(c)\right)</span>
<span class="l" id="L597"><span class="n">597</span> \left(\sum_{\substack{f^2\mid n\\(f,dR)=1}}\mu(f)\right).</span>
<span class="l" id="L598"><span class="n">598</span>\]</span>
<span class="l" id="L599"><span class="n">599</span>After writing $n=cf^2\ell$, every nonzero channel therefore has</span>
<span class="l" id="L600"><span class="n">600</span>\[</span>
<span class="l" id="L601"><span class="n">601</span> c\mid dR,\qquad (f,dR)=1,\qquad</span>
<span class="l" id="L602"><span class="n">602</span> (f,6q)=1,\qquad(\ell,6q)=1.</span>
<span class="l" id="L603"><span class="n">603</span>\]</span>
<span class="l" id="L604"><span class="n">604</span>In particular $d,c,f,\ell$ are units modulo $\Lambda_q$ wherever an inverse</span>
<span class="l" id="L605"><span class="n">605</span>modulo $\Lambda_q$ is evaluated.  If a moving reciprocal modulus is present,</span>
<span class="l" id="L606"><span class="n">606</span>its unit restriction is retained as a separate outside indicator through the</span>
<span class="l" id="L607"><span class="n">607</span>same split; it likewise forces every factor that is inverted modulo that moving</span>
<span class="l" id="L608"><span class="n">608</span>modulus to be a unit.  No inclusion--exclusion channel containing a factor $2$</span>
<span class="l" id="L609"><span class="n">609</span>or $3$ is ever asked to carry a $\Lambda_q$-reciprocal phase.</span>
<span class="l" id="L610"><span class="n">610</span></span>
<span class="l" id="L611"><span class="n">611</span>The next statement makes the fixed-conductor completion step explicit without</span>
<span class="l" id="L612"><span class="n">612</span>suppressing the zero Fourier mode of the reciprocal phase.  Let</span>
<span class="l" id="L613"><span class="n">613</span>$\Lambda=\Lambda_q$, let $r$ be a squarefree high modulus with</span>
<span class="l" id="L614"><span class="n">614</span>$(r,\Lambda)=1$, and let $I\subset(U,2U]$ be an interval.  For fixed integers</span>
<span class="l" id="L615"><span class="n">615</span>$A,B$ set</span>
<span class="l" id="L616"><span class="n">616</span>\begin{equation}\label{eq:fixed-conductor-sigma}</span>
<span class="l" id="L617"><span class="n">617</span>\Sigma_{r,\Lambda}^{B}(I;A,R)</span>
<span class="l" id="L618"><span class="n">618</span>=\sum_{\substack{\ell\in I\\(\ell,rR)=1}}</span>
<span class="l" id="L619"><span class="n">619</span> \omega_q(\ell)\,</span>
<span class="l" id="L620"><span class="n">620</span> \e_r(A\bar\ell)\,</span>
<span class="l" id="L621"><span class="n">621</span> \e_{\Lambda}(B\bar\ell).</span>
<span class="l" id="L622"><span class="n">622</span>\end{equation}</span>
<span class="l" id="L623"><span class="n">623</span>The support of $\omega_q$ already forces $(\ell,\Lambda)=1$.</span>
<span class="l" id="L624"><span class="n">624</span></span>
<span class="l" id="L625"><span class="n">625</span>\begin{lemma}[Centered fixed-conductor weighted completion]\label{lem:fixed-conductor-completion}</span>
<span class="l" id="L626"><span class="n">626</span>Put</span>
<span class="l" id="L627"><span class="n">627</span>\[</span>
<span class="l" id="L628"><span class="n">628</span>Q=r\Lambda,</span>
<span class="l" id="L629"><span class="n">629</span>\qquad C=A\Lambda+Br,</span>
<span class="l" id="L630"><span class="n">630</span>\qquad \rho=\frac{Q}{(C,Q)},</span>
<span class="l" id="L631"><span class="n">631</span>\qquad</span>
<span class="l" id="L632"><span class="n">632</span>\mathfrak m_Q(C)=\frac{\mu(\rho)}{\ph(\rho)},</span>
<span class="l" id="L633"><span class="n">633</span>\]</span>
<span class="l" id="L634"><span class="n">634</span>and</span>
<span class="l" id="L635"><span class="n">635</span>\[</span>
<span class="l" id="L636"><span class="n">636</span>M_{r,R}(I)=\sum_{\substack{\ell\in I\\(\ell,rR)=1}}\omega_q(\ell).</span>
<span class="l" id="L637"><span class="n">637</span>\]</span>
<span class="l" id="L638"><span class="n">638</span>Then</span>
<span class="l" id="L639"><span class="n">639</span>\begin{equation}\label{eq:combined-reciprocal-phase}</span>
<span class="l" id="L640"><span class="n">640</span>\e_r(A\bar\ell)\e_{\Lambda}(B\bar\ell)</span>
<span class="l" id="L641"><span class="n">641</span>=\e_Q(C\bar\ell)</span>
<span class="l" id="L642"><span class="n">642</span>\end{equation}</span>
<span class="l" id="L643"><span class="n">643</span>for every supported $\ell$, and</span>
<span class="l" id="L644"><span class="n">644</span>\begin{equation}\label{eq:fixed-conductor-completion}</span>
<span class="l" id="L645"><span class="n">645</span>\boxed{</span>
<span class="l" id="L646"><span class="n">646</span>\Sigma_{r,\Lambda}^{B}(I;A,R)</span>
<span class="l" id="L647"><span class="n">647</span>=\mathfrak m_Q(C)M_{r,R}(I)</span>
<span class="l" id="L648"><span class="n">648</span>+O_q\!\left(</span>
<span class="l" id="L649"><span class="n">649</span>\mathcal E(R)2^{\Om(r)/2}\rho^{1/4}U^{-1/2}(\log 2r)^{1/2}</span>
<span class="l" id="L650"><span class="n">650</span>\right),</span>
<span class="l" id="L651"><span class="n">651</span>}</span>
<span class="l" id="L652"><span class="n">652</span>\end{equation}</span>
<span class="l" id="L653"><span class="n">653</span>where</span>
<span class="l" id="L654"><span class="n">654</span>\[</span>
<span class="l" id="L655"><span class="n">655</span>\mathcal E(R)=\prod_{p\mid R}(1+p^{-1/2}),</span>
<span class="l" id="L656"><span class="n">656</span>\qquad M_{r,R}(I)\ll_q1.</span>
<span class="l" id="L657"><span class="n">657</span>\]</span>
<span class="l" id="L658"><span class="n">658</span>The implied constant is uniform in the endpoints of $I$.  Since $Q$ is</span>
<span class="l" id="L659"><span class="n">659</span>squarefree, the displayed main coefficient is the exact normalized Ramanujan</span>
<span class="l" id="L660"><span class="n">660</span>mean of the reciprocal phase.  Full resonance is exactly</span>
<span class="l" id="L661"><span class="n">661</span>\begin{equation}\label{eq:fixed-conductor-resonance}</span>
<span class="l" id="L662"><span class="n">662</span>\rho=1</span>
<span class="l" id="L663"><span class="n">663</span>\Longleftrightarrow</span>
<span class="l" id="L664"><span class="n">664</span>r\mid A\ \text{and}\ \Lambda\mid B.</span>
<span class="l" id="L665"><span class="n">665</span>\end{equation}</span>
<span class="l" id="L666"><span class="n">666</span>In particular, $\rho&gt;1$ does not by itself imply decay: the Ramanujan-mean</span>
<span class="l" id="L667"><span class="n">667</span>term in \eqref{eq:fixed-conductor-completion} must first be removed or summed.</span>
<span class="l" id="L668"><span class="n">668</span>\end{lemma}</span>
<span class="l" id="L669"><span class="n">669</span></span>
<span class="l" id="L670"><span class="n">670</span>\begin{proof}</span>
<span class="l" id="L671"><span class="n">671</span>Because $(r,\Lambda)=1$, CRT gives</span>
<span class="l" id="L672"><span class="n">672</span>\eqref{eq:combined-reciprocal-phase}.  The phase on the unit group modulo $Q$</span>
<span class="l" id="L673"><span class="n">673</span>has complete mean</span>
<span class="l" id="L674"><span class="n">674</span>\[</span>
<span class="l" id="L675"><span class="n">675</span>\frac1{\ph(Q)}\sum_{x\bmod Q}^{*}\e_Q(C\bar x)</span>
<span class="l" id="L676"><span class="n">676</span>=\frac{c_Q(C)}{\ph(Q)}</span>
<span class="l" id="L677"><span class="n">677</span>=\frac{\mu(\rho)}{\ph(\rho)}.</span>
<span class="l" id="L678"><span class="n">678</span>\]</span>
<span class="l" id="L679"><span class="n">679</span>Subtract this constant before completing an interval.  The centered periodic</span>
<span class="l" id="L680"><span class="n">680</span>function has zero sum on every complete $Q$-period, so only one incomplete</span>
<span class="l" id="L681"><span class="n">681</span>period remains.  Inclusion--exclusion over the resonant prime factors and</span>
<span class="l" id="L682"><span class="n">682</span>completion of that remainder on the reduced modulus $\rho$ give the same</span>
<span class="l" id="L683"><span class="n">683</span>squarefree Kloosterman bound as in the companion argument, now with the zero</span>
<span class="l" id="L684"><span class="n">684</span>Fourier mode absent.  Apply the unit-preserving split</span>
<span class="l" id="L685"><span class="n">685</span>\eqref{eq:fixed-q-squarefree-removal}: only divisors of the high-prime factor</span>
<span class="l" id="L686"><span class="n">686</span>$dR$ enter the $c$-sum, while $\1_{(n,6q)=1}$ stays outside the split.  Thus the</span>
<span class="l" id="L687"><span class="n">687</span>inner $\ell$-sum is always over units modulo $\Lambda$, and $c$ and $f$ are</span>
<span class="l" id="L688"><span class="n">688</span>units modulo $\Lambda$ as well.  Any moving reciprocal unit condition is kept</span>
<span class="l" id="L689"><span class="n">689</span>outside in the same way.  The centered incomplete Kloosterman estimate and the</span>
<span class="l" id="L690"><span class="n">690</span>finite $f$-balance are therefore exactly the companion estimates on each legal</span>
<span class="l" id="L691"><span class="n">691</span>unit channel, with at most an $O_q(1)$ fixed-period factor; this gives the error</span>
<span class="l" id="L692"><span class="n">692</span>term in \eqref{eq:fixed-conductor-completion}.  The untwisted dyadic mass</span>
<span class="l" id="L693"><span class="n">693</span>satisfies $M_{r,R}(I)\ll_q1$ by the elementary logarithmic mean for $\omega_q$,</span>
<span class="l" id="L694"><span class="n">694</span>and restrictions only delete nonnegative terms.  Finally, reducing</span>
<span class="l" id="L695"><span class="n">695</span>$C=A\Lambda+Br$ modulo $r$ and modulo $\Lambda$ proves</span>
<span class="l" id="L696"><span class="n">696</span>\eqref{eq:fixed-conductor-resonance}.</span>
<span class="l" id="L697"><span class="n">697</span>\end{proof}</span>
<span class="l" id="L698"><span class="n">698</span></span>
<span class="l" id="L699"><span class="n">699</span>Thus this fixed-$q$ port uses the companion squarefree-removal and finite</span>
<span class="l" id="L700"><span class="n">700</span>$f$-split architecture only after the reciprocal phase has been centered; it</span>
<span class="l" id="L701"><span class="n">701</span>does not invoke an arbitrary-interval estimate for the uncentered phase.</span>
<span class="l" id="L702"><span class="n">702</span></span>
<span class="l" id="L703"><span class="n">703</span>We use here exactly the companion exterior parameters</span>
<span class="l" id="L704"><span class="n">704</span>\[</span>
<span class="l" id="L705"><span class="n">705</span>L=\log h,</span>
<span class="l" id="L706"><span class="n">706</span>\qquad \sigma=\frac23+0.01,</span>
<span class="l" id="L707"><span class="n">707</span>\qquad K=L^{10},</span>
<span class="l" id="L708"><span class="n">708</span>\qquad W=\exp(L^\sigma),</span>
<span class="l" id="L709"><span class="n">709</span>\]</span>
<span class="l" id="L710"><span class="n">710</span>\[</span>
<span class="l" id="L711"><span class="n">711</span>T_{\rm comp}=\exp(L^{0.60}),</span>
<span class="l" id="L712"><span class="n">712</span>\qquad X=T_{\rm comp}^{3}W^2.</span>
<span class="l" id="L713"><span class="n">713</span>\]</span>
<span class="l" id="L714"><span class="n">714</span>We now record the mean-term ledger that is used in the exterior.  Write</span>
<span class="l" id="L715"><span class="n">715</span>\[</span>
<span class="l" id="L716"><span class="n">716</span>r_A=\frac r{(A,r)},</span>
<span class="l" id="L717"><span class="n">717</span>\qquad</span>
<span class="l" id="L718"><span class="n">718</span>\rho_\Lambda=\frac{\Lambda}{(B,\Lambda)}.</span>
<span class="l" id="L719"><span class="n">719</span>\]</span>
<span class="l" id="L720"><span class="n">720</span>Then $\rho=r_A\rho_\Lambda$ and hence</span>
<span class="l" id="L721"><span class="n">721</span>\begin{equation}\label{eq:fixed-conductor-mean-bound}</span>
<span class="l" id="L722"><span class="n">722</span>|\mathfrak m_Q(C)|\ll_q\frac1{\ph(r_A)}.</span>
<span class="l" id="L723"><span class="n">723</span>\end{equation}</span>
<span class="l" id="L724"><span class="n">724</span>In every actual nonzero Vaaler completion channel, the unit scalings introduced</span>
<span class="l" id="L725"><span class="n">725</span>by squarefree removal do not change the gcd with the moving modulus, and the</span>
<span class="l" id="L726"><span class="n">726</span>moving numerator is $jH$ times such a unit.  Since active moving moduli are</span>
<span class="l" id="L727"><span class="n">727</span>coprime to $H$ and $1\le |j|\le K$, one has</span>
<span class="l" id="L728"><span class="n">728</span>\[</span>
<span class="l" id="L729"><span class="n">729</span>r_A=\frac r{(r,j)}.</span>
<span class="l" id="L730"><span class="n">730</span>\]</span>
<span class="l" id="L731"><span class="n">731</span>For $Y\ge K$, Rankin&#x27;s trick at exponent $1/2$ gives, uniformly in</span>
<span class="l" id="L732"><span class="n">732</span>$1\le|j|\le K$,</span>
<span class="l" id="L733"><span class="n">733</span>\begin{align}</span>
<span class="l" id="L734"><span class="n">734</span>\sum_{r&gt;Y}\frac{\omega_q(r)}{\ph(r/(r,j))}</span>
<span class="l" id="L735"><span class="n">735</span>&amp;\ll_q</span>
<span class="l" id="L736"><span class="n">736</span>Y^{-1/2}</span>
<span class="l" id="L737"><span class="n">737</span>\sum_{d\mid\operatorname{rad}(j)}\omega_q(d)d^{1/2}\notag\\</span>
<span class="l" id="L738"><span class="n">738</span>&amp;\ll_q</span>
<span class="l" id="L739"><span class="n">739</span>Y^{-1/2}</span>
<span class="l" id="L740"><span class="n">740</span>\prod_{\substack{p\mid j\\p\ge5,\ p\nmid q}}</span>
<span class="l" id="L741"><span class="n">741</span>\left(1+\frac{\sqrt p}{p-3}\right)</span>
<span class="l" id="L742"><span class="n">742</span>=Y^{-1/2}K^{o(1)}.</span>
<span class="l" id="L743"><span class="n">743</span>\label{eq:ramanujan-mean-tail}</span>
<span class="l" id="L744"><span class="n">744</span>\end{align}</span>
<span class="l" id="L745"><span class="n">745</span>Indeed, after $r=ds$ with $d=(r,j)$, the remaining tail is bounded by</span>
<span class="l" id="L746"><span class="n">746</span>\[</span>
<span class="l" id="L747"><span class="n">747</span>\sum_{s&gt;Y/d}\frac{\omega_q(s)}{\ph(s)}</span>
<span class="l" id="L748"><span class="n">748</span>\le</span>
<span class="l" id="L749"><span class="n">749</span>\left(\frac dY\right)^{1/2}</span>
<span class="l" id="L750"><span class="n">750</span>\prod_{\substack{p\ge5\\p\nmid q}}</span>
<span class="l" id="L751"><span class="n">751</span>\left(1+\frac{\sqrt p}{(p-3)(p-1)}\right).</span>
<span class="l" id="L752"><span class="n">752</span>\]</span>
<span class="l" id="L753"><span class="n">753</span>The Euler product is convergent.  Since the other high-modulus mass up to</span>
<span class="l" id="L754"><span class="n">754</span>$h^2$ is $O_q(L)$, the main Vaaler frequencies have total $\ell^1$ mass</span>
<span class="l" id="L755"><span class="n">755</span>$O(\log K)$, and the fixed-period expansion has $O_q(1)$ terms, the aggregate</span>
<span class="l" id="L756"><span class="n">756</span>Ramanujan means from all moving moduli $r&gt;K$ contribute</span>
<span class="l" id="L757"><span class="n">757</span>\begin{equation}\label{eq:ramanujan-mean-global}</span>
<span class="l" id="L758"><span class="n">758</span>O_q\!\left(L\log K\,K^{-1/2+o(1)}\right)=o(1).</span>
<span class="l" id="L759"><span class="n">759</span>\end{equation}</span>
<span class="l" id="L760"><span class="n">760</span>For direct completion the moving modulus is in fact $&gt;X$, so</span>
<span class="l" id="L761"><span class="n">761</span>\eqref{eq:ramanujan-mean-tail} is stronger still.</span>
<span class="l" id="L762"><span class="n">762</span></span>
<span class="l" id="L763"><span class="n">763</span>No modewise completion is used for a moving modulus at most $K$.  In the</span>
<span class="l" id="L764"><span class="n">764</span>direct orientation we return, before expanding the fixed-period factor, to the</span>
<span class="l" id="L765"><span class="n">765</span>underlying sawtooth $\psi(x)=x-\lfloor x\rfloor-\tfrac12$ on the combined modulus $Q=r\Lambda$.  For every</span>
<span class="l" id="L766"><span class="n">766</span>nonconstant sawtooth channel its numerator $C_0$ is nonzero modulo $Q$, and</span>
<span class="l" id="L767"><span class="n">767</span>inversion permutes the units; pairing $x$ with $-x$ therefore gives the exact</span>
<span class="l" id="L768"><span class="n">768</span>complete-period identity</span>
<span class="l" id="L769"><span class="n">769</span>\begin{equation}\label{eq:fixed-conductor-sawtooth-zero}</span>
<span class="l" id="L770"><span class="n">770</span>\sum_{x\bmod Q}^{*}\psi(C_0\bar x/Q)=0.</span>
<span class="l" id="L771"><span class="n">771</span>\end{equation}</span>
<span class="l" id="L772"><span class="n">772</span>The constant channel is the already separated deterministic/full-resonance</span>
<span class="l" id="L773"><span class="n">773</span>piece.  The finite $f$-split applied to the remaining incomplete $Q$-period</span>
<span class="l" id="L774"><span class="n">774</span>then gives, with only a fixed-$q$ change in the constant,</span>
<span class="l" id="L775"><span class="n">775</span>\begin{equation}\label{eq:fixed-conductor-small-modulus}</span>
<span class="l" id="L776"><span class="n">776</span>\ll_q</span>
<span class="l" id="L777"><span class="n">777</span>\mathcal E(R)(\log 2K)^{C_1}(K/U)^{1/2},</span>
<span class="l" id="L778"><span class="n">778</span>\end{equation}</span>
<span class="l" id="L779"><span class="n">779</span>for a fixed absolute $C_1$.  This is super-logarithmically small in the exterior range $U&gt;X$.  Thus the</span>
<span class="l" id="L780"><span class="n">780</span>nonzero Fourier modes are never charged individually for their nonzero</span>
<span class="l" id="L781"><span class="n">781</span>Ramanujan means when the moving modulus is small.</span>
<span class="l" id="L782"><span class="n">782</span></span>
<span class="l" id="L783"><span class="n">783</span>In the reciprocal orientation the companion geometric carve-out already</span>
<span class="l" id="L784"><span class="n">784</span>removes the whole region with moving modulus $u\le K$ before slow/moderate</span>
<span class="l" id="L785"><span class="n">785</span>completion.  We retain that whole carve-out for every fixed-period frequency,</span>
<span class="l" id="L786"><span class="n">786</span>not only for the principal one.  Its band and tail remain respectively</span>
<span class="l" id="L787"><span class="n">787</span>\[</span>
<span class="l" id="L788"><span class="n">788</span>O_q(L^\sigma\log L)=o(L),</span>
<span class="l" id="L789"><span class="n">789</span>\qquad O_q(1),</span>
<span class="l" id="L790"><span class="n">790</span>\]</span>
<span class="l" id="L791"><span class="n">791</span>because the fixed-conductor factors are bounded.  Outside the carve-out the</span>
<span class="l" id="L792"><span class="n">792</span>centered error in \eqref{eq:fixed-conductor-completion} passes through the same</span>
<span class="l" id="L793"><span class="n">793</span>weighted-completion and reciprocal estimates as before; the fixed factor</span>
<span class="l" id="L794"><span class="n">794</span>changes only $O_q(1)$ constants.  The fast reciprocal channels are unchanged:</span>
<span class="l" id="L795"><span class="n">795</span>finite Fourier expansion of the fixed-period decoration adds only linear prime</span>
<span class="l" id="L796"><span class="n">796</span>phases, and Karatsuba&#x27;s derivative tests used here \cite[Theorem~1]{Karatsuba1971} begin at order at least two.</span>
<span class="l" id="L797"><span class="n">797</span>The Rankin tail and inherited forbidden-prime surcharge likewise change only</span>
<span class="l" id="L798"><span class="n">798</span>by fixed-$q$ constants.</span>
<span class="l" id="L799"><span class="n">799</span></span>
<span class="l" id="L800"><span class="n">800</span>Combining the centered completion error, the exact small-modulus sawtooth</span>
<span class="l" id="L801"><span class="n">801</span>cancellation, the retained reciprocal carve-out, and the negligible global</span>
<span class="l" id="L802"><span class="n">802</span>Ramanujan-mean ledger \eqref{eq:ramanujan-mean-global}, the complete hyperbolic</span>
<span class="l" id="L803"><span class="n">803</span>exterior, including the Vaaler remainder frequencies and the slow</span>
<span class="l" id="L804"><span class="n">804</span>moving-resonance contribution, is</span>
<span class="l" id="L805"><span class="n">805</span>\begin{equation}\label{eq:exterior-q}</span>
<span class="l" id="L806"><span class="n">806</span>\boxed{o_q(\log h)}</span>
<span class="l" id="L807"><span class="n">807</span>\end{equation}</span>
<span class="l" id="L808"><span class="n">808</span>pointwise.</span>
<span class="l" id="L809"><span class="n">809</span></span>
<span class="l" id="L810"><span class="n">810</span>\subsection{Decorated rows}</span>
<span class="l" id="L811"><span class="n">811</span></span>
<span class="l" id="L812"><span class="n">812</span>These rows carry a deterministic divisor of the height.  For a height-divisor</span>
<span class="l" id="L813"><span class="n">813</span>row write $e_1=dr$ with $d\mid h$ and $(d,r)=1$.  All active high primes are</span>
<span class="l" id="L814"><span class="n">814</span>coprime to $q$, so $(d,q)=1$.  Under the exact change of variables $y=dz$,</span>
<span class="l" id="L815"><span class="n">815</span>\begin{equation}\label{eq:decorated-residue}</span>
<span class="l" id="L816"><span class="n">816</span>(a+y,q)=1</span>
<span class="l" id="L817"><span class="n">817</span>\Longleftrightarrow</span>
<span class="l" id="L818"><span class="n">818</span>(d^{-1}a+z,q)=1.</span>
<span class="l" id="L819"><span class="n">819</span>\end{equation}</span>
<span class="l" id="L820"><span class="n">820</span>The low $2,3$ ratio is invariant when both the height and the free coordinate</span>
<span class="l" id="L821"><span class="n">821</span>are multiplied by a high unit modulo $6$.  Therefore a decorated row makes only</span>
<span class="l" id="L822"><span class="n">822</span>the finite residue permutation</span>
<span class="l" id="L823"><span class="n">823</span>\begin{equation}\label{eq:decorated-map}</span>
<span class="l" id="L824"><span class="n">824</span>(a,h)\longmapsto(d^{-1}a,h/d)</span>
<span class="l" id="L825"><span class="n">825</span>\end{equation}</span>
<span class="l" id="L826"><span class="n">826</span>inside the fixed family of endpoint classes.</span>
<span class="l" id="L827"><span class="n">827</span></span>
<span class="l" id="L828"><span class="n">828</span>The exterior estimate at this reduced height is run with the centered</span>
<span class="l" id="L829"><span class="n">829</span>completion just proved, not with an uncentered companion estimate.  The actual</span>
<span class="l" id="L830"><span class="n">830</span>row height changes only the factor $H$ in the moving numerator.  Active moving</span>
<span class="l" id="L831"><span class="n">831</span>moduli remain coprime to that height, so in every Vaaler channel</span>
<span class="l" id="L832"><span class="n">832</span>$r_A=r/(r,j)$ exactly as in \eqref{eq:ramanujan-mean-tail}.  Hence the direct</span>
<span class="l" id="L833"><span class="n">833</span>orientation uses \eqref{eq:fixed-conductor-completion} together with the mean</span>
<span class="l" id="L834"><span class="n">834</span>ledger \eqref{eq:ramanujan-mean-global}; the reciprocal orientation retains the</span>
<span class="l" id="L835"><span class="n">835</span>whole moving-modulus-$\le K$ carve-out and uses the same centered estimate and</span>
<span class="l" id="L836"><span class="n">836</span>mean ledger outside it.  Both statements are uniform for actual row heights at</span>
<span class="l" id="L837"><span class="n">837</span>most the ambient $h$.  The outer decoration norm is $O(\log\log h)$, while</span>
<span class="l" id="L838"><span class="n">838</span>\eqref{eq:ramanujan-mean-global} is $O_q(L\log K\,K^{-1/2+o(1)})=o_q(1)$ and</span>
<span class="l" id="L839"><span class="n">839</span>the centered nonzero-frequency saving is arbitrarily logarithmic.  Therefore</span>
<span class="l" id="L840"><span class="n">840</span>summing all decorated rows still gives $o_q(\log h)$.</span>
<span class="l" id="L841"><span class="n">841</span></span>
<span class="l" id="L842"><span class="n">842</span>\subsection{One-dimensional boundary terms}</span>
<span class="l" id="L843"><span class="n">843</span></span>
<span class="l" id="L844"><span class="n">844</span>The two boundary families are matched before they are estimated.  Exact reflection turns the right boundary into a left boundary with residue $-b$; at divisor-family level, the same reflection leaves only the subfamily supported entirely on active primes dividing $h$, which is smaller by an $O_q(\log\log h)$ estimate.  Define the left free one-anchor sum</span>
<span class="l" id="L845"><span class="n">845</span>\begin{equation}\label{eq:Sqa}</span>
<span class="l" id="L846"><span class="n">846</span>S_{q,a}(h)=\sum_{t\in I_q(a;h)}(\SSer_q(\{0,t\})-1).</span>
<span class="l" id="L847"><span class="n">847</span>\end{equation}</span>
<span class="l" id="L848"><span class="n">848</span>Reflection is exact.  If $u=h-t$, then</span>
<span class="l" id="L849"><span class="n">849</span>\[</span>
<span class="l" id="L850"><span class="n">850</span>(a+t,q)=1</span>
<span class="l" id="L851"><span class="n">851</span>\Longleftrightarrow</span>
<span class="l" id="L852"><span class="n">852</span>(b-u,q)=1</span>
<span class="l" id="L853"><span class="n">853</span>\Longleftrightarrow</span>
<span class="l" id="L854"><span class="n">854</span>(-b+u,q)=1,</span>
<span class="l" id="L855"><span class="n">855</span>\]</span>
<span class="l" id="L856"><span class="n">856</span>and $\SSer_q(\{h,t\})=\SSer_q(\{0,u\})$.  Hence</span>
<span class="l" id="L857"><span class="n">857</span>\begin{equation}\label{eq:reflection-q}</span>
<span class="l" id="L858"><span class="n">858</span>\boxed{</span>
<span class="l" id="L859"><span class="n">859</span>\sum_{t\in I_q(a;h)}(\SSer_q(\{h,t\})-1)</span>
<span class="l" id="L860"><span class="n">860</span>=S_{q,-b}(h).</span>
<span class="l" id="L861"><span class="n">861</span>}</span>
<span class="l" id="L862"><span class="n">862</span>\end{equation}</span>
<span class="l" id="L863"><span class="n">863</span></span>
<span class="l" id="L864"><span class="n">864</span>The divisor families require the corresponding row-level reflection, not only</span>
<span class="l" id="L865"><span class="n">865</span>this free-boundary identity.  For a unit residue $r\pmod q$, write</span>
<span class="l" id="L866"><span class="n">866</span>$\lambda_{q,r,h}$, $N_q(r;h)$, and $\Delta_{q,r}$ for</span>
<span class="l" id="L867"><span class="n">867</span>\eqref{eq:lambdaq}, \eqref{eq:IqNq}, and \eqref{eq:Deltaqea} with</span>
<span class="l" id="L868"><span class="n">868</span>$a$ replaced by $r$, and put</span>
<span class="l" id="L869"><span class="n">869</span>\[</span>
<span class="l" id="L870"><span class="n">870</span>R_{q,h}:=\prod_{\substack{p\mid h\\p\ge5,\ p\nmid q}}p.</span>
<span class="l" id="L871"><span class="n">871</span>\]</span>
<span class="l" id="L872"><span class="n">872</span>Define the residue-sensitive pure-one family</span>
<span class="l" id="L873"><span class="n">873</span>\begin{equation}\label{eq:P1-residue-q}</span>
<span class="l" id="L874"><span class="n">874</span>P_{1,q}^{(r)}(h)</span>
<span class="l" id="L875"><span class="n">875</span>:=C_q^*(h)\sum_{f&gt;1}v_{q,h}(f)\Delta_{q,r}(h;f,1),</span>
<span class="l" id="L876"><span class="n">876</span>\end{equation}</span>
<span class="l" id="L877"><span class="n">877</span>and its pure-height-prime support subfamily</span>
<span class="l" id="L878"><span class="n">878</span>\begin{equation}\label{eq:P1-support-q}</span>
<span class="l" id="L879"><span class="n">879</span>P_{1,q,h}^{\mathrm{supp}}(r)</span>
<span class="l" id="L880"><span class="n">880</span>:=C_q^*(h)</span>
<span class="l" id="L881"><span class="n">881</span>\sum_{\substack{f&gt;1\\f\mid R_{q,h}}}</span>
<span class="l" id="L882"><span class="n">882</span>v_{q,h}(f)\Delta_{q,r}(h;f,1).</span>
<span class="l" id="L883"><span class="n">883</span>\end{equation}</span>
<span class="l" id="L884"><span class="n">884</span>Here $P_{2,q}$ denotes the contribution to \eqref{eq:Tq-expansion} with</span>
<span class="l" id="L885"><span class="n">885</span>$e_1=1&lt;e_2$, and $D_{1,q}$ denotes the contribution with</span>
<span class="l" id="L886"><span class="n">886</span>$e_1=d&gt;1$, $d\mid R_{q,h}$, $e_2&gt;1$, and $(e_2,h)=1$.</span>
<span class="l" id="L887"><span class="n">887</span></span>
<span class="l" id="L888"><span class="n">888</span>\begin{lemma}[Fixed-$q$ divisor-family reflection]\label{lem:reflection-family-q}</span>
<span class="l" id="L889"><span class="n">889</span>Let $d\mid R_{q,h}$ and let $e$ lie on the support of $\omega_q$ with</span>
<span class="l" id="L890"><span class="n">890</span>$(e,h)=1$.  Then</span>
<span class="l" id="L891"><span class="n">891</span>\begin{equation}\label{eq:reflection-row-q}</span>
<span class="l" id="L892"><span class="n">892</span>\boxed{\Delta_{q,a}(h;d,e)=\Delta_{q,-b}(h;de,1),}</span>
<span class="l" id="L893"><span class="n">893</span>\end{equation}</span>
<span class="l" id="L894"><span class="n">894</span>and</span>
<span class="l" id="L895"><span class="n">895</span>\begin{equation}\label{eq:reflection-weight-q}</span>
<span class="l" id="L896"><span class="n">896</span>\boxed{v_{q,h}(d)\omega_q(e)=v_{q,h}(de).}</span>
<span class="l" id="L897"><span class="n">897</span>\end{equation}</span>
<span class="l" id="L898"><span class="n">898</span>Consequently the one-dimensional divisor families satisfy the exact identity</span>
<span class="l" id="L899"><span class="n">899</span>\begin{equation}\label{eq:reflection-assembly-q}</span>
<span class="l" id="L900"><span class="n">900</span>\boxed{</span>
<span class="l" id="L901"><span class="n">901</span>P_{2,q}+D_{1,q}</span>
<span class="l" id="L902"><span class="n">902</span>=P_{1,q}^{(-b)}-P_{1,q,h}^{\mathrm{supp}}(-b),</span>
<span class="l" id="L903"><span class="n">903</span>}</span>
<span class="l" id="L904"><span class="n">904</span>\end{equation}</span>
<span class="l" id="L905"><span class="n">905</span>with</span>
<span class="l" id="L906"><span class="n">906</span>\begin{equation}\label{eq:support-bound-q}</span>
<span class="l" id="L907"><span class="n">907</span>P_{1,q,h}^{\mathrm{supp}}(r)=O_q(\log\log h)</span>
<span class="l" id="L908"><span class="n">908</span>\end{equation}</span>
<span class="l" id="L909"><span class="n">909</span>uniformly over the finitely many compatible unit residues $r\pmod q$.</span>
<span class="l" id="L910"><span class="n">910</span>\end{lemma}</span>
<span class="l" id="L911"><span class="n">911</span></span>
<span class="l" id="L912"><span class="n">912</span>\begin{proof}</span>
<span class="l" id="L913"><span class="n">913</span>Put $z=h-y$.  Reflection is an affine bijection modulo every active low</span>
<span class="l" id="L914"><span class="n">914</span>prime, so</span>
<span class="l" id="L915"><span class="n">915</span>$L_p(\{0,y,h\})=L_p(\{0,z,h\})$ for every active</span>
<span class="l" id="L916"><span class="n">916</span>$p\in\{2,3\}$.</span>
<span class="l" id="L917"><span class="n">917</span>Compatibility $h\equiv b-a\pmod q$ gives</span>
<span class="l" id="L918"><span class="n">918</span>\[</span>
<span class="l" id="L919"><span class="n">919</span>(a+y,q)=1</span>
<span class="l" id="L920"><span class="n">920</span>\Longleftrightarrow (b-z,q)=1</span>
<span class="l" id="L921"><span class="n">921</span>\Longleftrightarrow (-b+z,q)=1.</span>
<span class="l" id="L922"><span class="n">922</span>\]</span>
<span class="l" id="L923"><span class="n">923</span>Hence</span>
<span class="l" id="L924"><span class="n">924</span>$\lambda_{q,a,h}(y)=\lambda_{q,-b,h}(z)$.</span>
<span class="l" id="L925"><span class="n">925</span>The same reflection bijects $I_q(a;h)$ with $I_q(-b;h)$, so</span>
<span class="l" id="L926"><span class="n">926</span>$N_q(a;h)=N_q(-b;h)$ exactly.  If $d\mid h$, then</span>
<span class="l" id="L927"><span class="n">927</span>$d\mid y$ if and only if $d\mid z$; the condition $e\mid h-y$</span>
<span class="l" id="L928"><span class="n">928</span>is exactly $e\mid z$.  Since $(d,e)=1$, these are together equivalent</span>
<span class="l" id="L929"><span class="n">929</span>to $de\mid z$.  Subtracting the identical exact density denominator</span>
<span class="l" id="L930"><span class="n">930</span>therefore proves \eqref{eq:reflection-row-q}.  On the coefficient side,</span>
<span class="l" id="L931"><span class="n">931</span>all primes of $d$ divide $h$ while all primes of $e$ do not, so the local</span>
<span class="l" id="L932"><span class="n">932</span>definitions \eqref{eq:omegaq} and \eqref{eq:vqh} give</span>
<span class="l" id="L933"><span class="n">933</span>\eqref{eq:reflection-weight-q} prime by prime.</span>
<span class="l" id="L934"><span class="n">934</span></span>
<span class="l" id="L935"><span class="n">935</span>The family $P_{2,q}$ is the case $d=1$, $e&gt;1$; $D_{1,q}$ is the case</span>
<span class="l" id="L936"><span class="n">936</span>$d&gt;1$, $e&gt;1$.  After reflection these are exactly the pure-one indices</span>
<span class="l" id="L937"><span class="n">937</span>$f=de&gt;1$ having at least one active high prime not dividing $h$.  The</span>
<span class="l" id="L938"><span class="n">938</span>indices missing from $P_{1,q}^{(-b)}$ are therefore precisely the</span>
<span class="l" id="L939"><span class="n">939</span>$f&gt;1$ supported on $R_{q,h}$, which proves</span>
<span class="l" id="L940"><span class="n">940</span>\eqref{eq:reflection-assembly-q}.</span>
<span class="l" id="L941"><span class="n">941</span></span>
<span class="l" id="L942"><span class="n">942</span>It remains to estimate this support family.  If $f\mid R_{q,h}$ and $H=h/f$,</span>
<span class="l" id="L943"><span class="n">943</span>then $f$ is a unit modulo $\Lambda_q$.  Along the progression $y=fx$,</span>
<span class="l" id="L944"><span class="n">944</span>the low/admissibility weight is a bounded $\Lambda_q$-periodic sequence</span>
<span class="l" id="L945"><span class="n">945</span>whose complete-period mean is $\ph(q)/q$: multiplication by $f$ merely</span>
<span class="l" id="L946"><span class="n">946</span>permutes the residue classes modulo $q$, $2$, and $3$, while the active</span>
<span class="l" id="L947"><span class="n">947</span>low-prime local tables have mean one.  Thus</span>
<span class="l" id="L948"><span class="n">948</span>\[</span>
<span class="l" id="L949"><span class="n">949</span>N_{q,r}(h;f,1)=\frac{\ph(q)}q H+O_q(1).</span>
<span class="l" id="L950"><span class="n">950</span>\]</span>
<span class="l" id="L951"><span class="n">951</span>By the exact periodic count \eqref{eq:Nq-periodic},</span>
<span class="l" id="L952"><span class="n">952</span>$N_q(r;h)/f=(\ph(q)/q)H+O_q(1)$, uniformly in $f$, and hence</span>
<span class="l" id="L953"><span class="n">953</span>\[</span>
<span class="l" id="L954"><span class="n">954</span>\Delta_{q,r}(h;f,1)=O_q(1)</span>
<span class="l" id="L955"><span class="n">955</span>\qquad(f\mid R_{q,h}).</span>
<span class="l" id="L956"><span class="n">956</span>\]</span>
<span class="l" id="L957"><span class="n">957</span>The product $C_q^*(h)$ is bounded for fixed $q$, uniformly in $h$, and on</span>
<span class="l" id="L958"><span class="n">958</span>this support</span>
<span class="l" id="L959"><span class="n">959</span>$v_{q,h}(f)=\prod_{p\mid f}(p-2)^{-1}$.  Therefore</span>
<span class="l" id="L960"><span class="n">960</span>\[</span>
<span class="l" id="L961"><span class="n">961</span>|P_{1,q,h}^{\mathrm{supp}}(r)|</span>
<span class="l" id="L962"><span class="n">962</span>\ll_q</span>
<span class="l" id="L963"><span class="n">963</span>\sum_{f\mid R_{q,h}}v_{q,h}(f)</span>
<span class="l" id="L964"><span class="n">964</span>=\prod_{p\mid R_{q,h}}\left(1+\frac1{p-2}\right)</span>
<span class="l" id="L965"><span class="n">965</span>\ll_q\log\log h,</span>
<span class="l" id="L966"><span class="n">966</span>\]</span>
<span class="l" id="L967"><span class="n">967</span>by the standard Mertens bound, proving \eqref{eq:support-bound-q}.</span>
<span class="l" id="L968"><span class="n">968</span>\end{proof}</span>
<span class="l" id="L969"><span class="n">969</span></span>
<span class="l" id="L970"><span class="n">970</span>For a one-dimensional row let $H$ again be the actual parity-reduced kernel</span>
<span class="l" id="L971"><span class="n">971</span>height and set</span>
<span class="l" id="L972"><span class="n">972</span>\begin{equation}\label{eq:X0-contract}</span>
<span class="l" id="L973"><span class="n">973</span>\boxed{X_0=H-1.}</span>
<span class="l" id="L974"><span class="n">974</span>\end{equation}</span>
<span class="l" id="L975"><span class="n">975</span>The symbol $X$ below is only the current upper limit in an $e$-sum,</span>
<span class="l" id="L976"><span class="n">976</span>$1\le X\le X_0$; it is not a second height.  In the P\&#x27;etermann split put</span>
<span class="l" id="L977"><span class="n">977</span>\[</span>
<span class="l" id="L978"><span class="n">978</span> z=X_0/\Lambda_q,</span>
<span class="l" id="L979"><span class="n">979</span> \qquad</span>
<span class="l" id="L980"><span class="n">980</span> Y=z\exp\!\left(- (\log z)^b\right).</span>
<span class="l" id="L981"><span class="n">981</span>\]</span>
<span class="l" id="L982"><span class="n">982</span>The coefficient $v(e)$ in any row is explicitly one of $\omega_q(e)$,</span>
<span class="l" id="L983"><span class="n">983</span>$\eta_q(e)$, or $v_{q,h}(e)$ from</span>
<span class="l" id="L984"><span class="n">984</span>\eqref{eq:omegaq}, \eqref{eq:etaq}, and \eqref{eq:vqh}.</span>
<span class="l" id="L985"><span class="n">985</span></span>
<span class="l" id="L986"><span class="n">986</span>After splitting a one-dimensional root count by $c\bmod\Lambda_q$ and then</span>
<span class="l" id="L987"><span class="n">987</span>by a unit class modulo $\Lambda_q$, character orthogonality reduces every</span>
<span class="l" id="L988"><span class="n">988</span>nonprincipal periodic contribution to a finite linear combination of sums</span>
<span class="l" id="L989"><span class="n">989</span>\begin{equation}\label{eq:petermann-periodic}</span>
<span class="l" id="L990"><span class="n">990</span>\sum_{e\le X}\chi(e)v(e)</span>
<span class="l" id="L991"><span class="n">991</span>\psi\!\left(\frac{X_0}{\Lambda_q e}+\frac r{\Lambda_q}\right),</span>
<span class="l" id="L992"><span class="n">992</span>\end{equation}</span>
<span class="l" id="L993"><span class="n">993</span>where $\chi$ is a fixed character modulo $\Lambda_q$, $r$ is fixed, and</span>
<span class="l" id="L994"><span class="n">994</span>\[\psi(x)=x-\lfloor x\rfloor-\tfrac12.\]</span>
<span class="l" id="L995"><span class="n">995</span></span>
<span class="l" id="L996"><span class="n">996</span>\begin{lemma}[Fixed-conductor character/rational-shift P\&#x27;etermann estimate]\label{lem:fixed-conductor-petermann}</span>
<span class="l" id="L997"><span class="n">997</span>Fix $\Lambda\ge1$.  Let $v$ be one of the squarefree nonnegative base</span>
<span class="l" id="L998"><span class="n">998</span>coefficients occurring in the one-dimensional argument after deleting finitely</span>
<span class="l" id="L999"><span class="n">999</span>many primes, and suppose the ordinary $\alpha=0$ P\&#x27;etermann hypotheses hold.</span>
<span class="l" id="L1000"><span class="n">1000</span>For a fixed Dirichlet character $\chi\pmod\Lambda$, fixed integer $r$,</span>
<span class="l" id="L1001"><span class="n">1001</span>$0&lt;b&lt;1$, and a real reciprocal numerator $T\ge2\Lambda$, put</span>
<span class="l" id="L1002"><span class="n">1002</span>\[</span>
<span class="l" id="L1003"><span class="n">1003</span> Y_T=\frac{T}{\Lambda}</span>
<span class="l" id="L1004"><span class="n">1004</span> \exp\!\left(-\bigl(\log(T/\Lambda)\bigr)^b\right).</span>
<span class="l" id="L1005"><span class="n">1005</span>\]</span>
<span class="l" id="L1006"><span class="n">1006</span>Then</span>
<span class="l" id="L1007"><span class="n">1007</span>\[</span>
<span class="l" id="L1008"><span class="n">1008</span> \sum_{e\le Y_T}\chi(e)v(e)</span>
<span class="l" id="L1009"><span class="n">1009</span> \psi\!\left(\frac{T}{\Lambda e}+\frac r\Lambda\right)</span>
<span class="l" id="L1010"><span class="n">1010</span> \ll_{\Lambda,b}</span>
<span class="l" id="L1011"><span class="n">1011</span> (\log T)^{2/3}(\log\log(3T))^{4/3}+\log\log(3T).</span>
<span class="l" id="L1012"><span class="n">1012</span>\]</span>
<span class="l" id="L1013"><span class="n">1013</span>For $b=3/4$ the same endpoint splitting as in the ordinary P\&#x27;etermann argument</span>
<span class="l" id="L1014"><span class="n">1014</span>gives, uniformly for the partial sums used below,</span>
<span class="l" id="L1015"><span class="n">1015</span>\begin{equation}\label{eq:petermann-real-numerator}</span>
<span class="l" id="L1016"><span class="n">1016</span> \sum_{e\le X}\chi(e)v(e)</span>
<span class="l" id="L1017"><span class="n">1017</span> \psi\!\left(\frac{T}{\Lambda e}+\frac r\Lambda\right)</span>
<span class="l" id="L1018"><span class="n">1018</span> \ll_{\Lambda}</span>
<span class="l" id="L1019"><span class="n">1019</span> (\log(3T))^{3/4}</span>
<span class="l" id="L1020"><span class="n">1020</span> +(\log(3T))^{2/3}(\log\log(3T))^{4/3}</span>
<span class="l" id="L1021"><span class="n">1021</span> +\log\log(3T)+1</span>
<span class="l" id="L1022"><span class="n">1022</span>\end{equation}</span>
<span class="l" id="L1023"><span class="n">1023</span>for $1\le X\le T$.  If $T&lt;2\Lambda$, the corresponding sum has</span>
<span class="l" id="L1024"><span class="n">1024</span>$O_\Lambda(1)$ terms and is $O_\Lambda(1)$.</span>
<span class="l" id="L1025"><span class="n">1025</span>\end{lemma}</span>
<span class="l" id="L1026"><span class="n">1026</span></span>
<span class="l" id="L1027"><span class="n">1027</span>\begin{proof}</span>
<span class="l" id="L1028"><span class="n">1028</span>On the supported unit integers, expand the fixed character by finite additive</span>
<span class="l" id="L1029"><span class="n">1029</span>Fourier inversion,</span>
<span class="l" id="L1030"><span class="n">1030</span>\[</span>
<span class="l" id="L1031"><span class="n">1031</span> \chi(n)=\sum_{a\bmod\Lambda}\widehat\chi(a)\e(an/\Lambda),</span>
<span class="l" id="L1032"><span class="n">1032</span> \qquad \sum_a|\widehat\chi(a)|=O_\Lambda(1).</span>
<span class="l" id="L1033"><span class="n">1033</span>\]</span>
<span class="l" id="L1034"><span class="n">1034</span>Thus each prime phase in the P\&#x27;etermann exponential-sum stage is changed from</span>
<span class="l" id="L1035"><span class="n">1035</span>$F(p)$ to $F(p)+ap/\Lambda$.  Every derivative test used there starts at order</span>
<span class="l" id="L1036"><span class="n">1036</span>at least two, so all such derivative hypotheses are unchanged.  The monotone</span>
<span class="l" id="L1037"><span class="n">1037</span>positive coefficient $pv(p)$ remains outside Abel summation and the character</span>
<span class="l" id="L1038"><span class="n">1038</span>remains in the already controlled oscillatory prime partial sum; prime powers</span>
<span class="l" id="L1039"><span class="n">1039</span>$\nu\ge2$ vanish by squarefree support.  Absolute and Cauchy steps are</span>
<span class="l" id="L1040"><span class="n">1040</span>unchanged.  The reciprocal parameter $T$ is allowed to be real throughout</span>
<span class="l" id="L1041"><span class="n">1041</span>those derivative estimates.</span>
<span class="l" id="L1042"><span class="n">1042</span></span>
<span class="l" id="L1043"><span class="n">1043</span>The rational shift multiplies every nonzero sawtooth Fourier mode by the unit</span>
<span class="l" id="L1044"><span class="n">1044</span>factor $\e(jr/\Lambda)$ and rescales the reciprocal parameter by the fixed</span>
<span class="l" id="L1045"><span class="n">1045</span>factor $1/\Lambda$.  With smoothing width $1/T$, a jump crossing for</span>
<span class="l" id="L1046"><span class="n">1046</span>$e&lt;T/\Lambda$ would require</span>
<span class="l" id="L1047"><span class="n">1047</span>\[</span>
<span class="l" id="L1048"><span class="n">1048</span> T+re&lt;\Lambda e n\le T+re+\frac{\Lambda e}{T}.</span>
<span class="l" id="L1049"><span class="n">1049</span>\]</span>
<span class="l" id="L1050"><span class="n">1050</span>The last interval has length less than one.  Since</span>
<span class="l" id="L1051"><span class="n">1051</span>$\Lambda e n-re=e(\Lambda n-r)$ is an integer, any crossing must satisfy</span>
<span class="l" id="L1052"><span class="n">1052</span>\[</span>
<span class="l" id="L1053"><span class="n">1053</span> e(\Lambda n-r)=\lceil T\rceil.</span>
<span class="l" id="L1054"><span class="n">1054</span>\]</span>
<span class="l" id="L1055"><span class="n">1055</span>Thus every crossing index $e$ divides $\lceil T\rceil$ (and obeys one fixed</span>
<span class="l" id="L1056"><span class="n">1056</span>congruence for the complementary factor).  For either base coefficient used</span>
<span class="l" id="L1057"><span class="n">1057</span>here,</span>
<span class="l" id="L1058"><span class="n">1058</span>\[</span>
<span class="l" id="L1059"><span class="n">1059</span> \sum_{e\mid\lceil T\rceil}v(e)</span>
<span class="l" id="L1060"><span class="n">1060</span> \ll_\Lambda</span>
<span class="l" id="L1061"><span class="n">1061</span> \prod_{\substack{p\mid\lceil T\rceil\\p\ge5}}</span>
<span class="l" id="L1062"><span class="n">1062</span> \left(1+\frac1{p-3}\right)</span>
<span class="l" id="L1063"><span class="n">1063</span> \ll_\Lambda\log\log(3T),</span>
<span class="l" id="L1064"><span class="n">1064</span>\]</span>
<span class="l" id="L1065"><span class="n">1065</span>with the same or smaller bound for the $1/(p-2)$ local factor.  This is the</span>
<span class="l" id="L1066"><span class="n">1066</span>entire crossing bill.  The complementary range $T/\Lambda\le e\le T$ has</span>
<span class="l" id="L1067"><span class="n">1067</span>$O_\Lambda(1)$ absolute coefficient mass by the ordinary mean formula, and the</span>
<span class="l" id="L1068"><span class="n">1068</span>usual strip from $Y_T$ to $T/\Lambda$ has mass</span>
<span class="l" id="L1069"><span class="n">1069</span>$O_\Lambda((\log T)^b)$.  Adding the divisor-supported crossing bill proves the</span>
<span class="l" id="L1070"><span class="n">1070</span>claims.</span>
<span class="l" id="L1071"><span class="n">1071</span>\end{proof}</span>
<span class="l" id="L1072"><span class="n">1072</span></span>
<span class="l" id="L1073"><span class="n">1073</span>For the moving coefficient, the convolution used below is local and exact.</span>
<span class="l" id="L1074"><span class="n">1074</span>At each active high prime $p\mid h$ put</span>
<span class="l" id="L1075"><span class="n">1075</span>\[</span>
<span class="l" id="L1076"><span class="n">1076</span> a_p=\frac1{p-3},\qquad b_p=\frac1{p-2},</span>
<span class="l" id="L1077"><span class="n">1077</span>\]</span>
<span class="l" id="L1078"><span class="n">1078</span>and define the multiplicative correction $c_{q,h}$ by the Euler factors</span>
<span class="l" id="L1079"><span class="n">1079</span>\begin{equation}\label{eq:cqh-definition}</span>
<span class="l" id="L1080"><span class="n">1080</span>\frac{1+b_pz}{1+a_pz}</span>
<span class="l" id="L1081"><span class="n">1081</span>=1+\sum_{\nu\ge1}c_{q,h}(p^\nu)z^\nu,</span>
<span class="l" id="L1082"><span class="n">1082</span>\qquad</span>
<span class="l" id="L1083"><span class="n">1083</span>c_{q,h}(p^\nu)=(b_p-a_p)(-a_p)^{\nu-1},</span>
<span class="l" id="L1084"><span class="n">1084</span>\end{equation}</span>
<span class="l" id="L1085"><span class="n">1085</span>at primes $p\ge5$ with $p\nmid q$ and $p\mid h$, and by the factor $1$</span>
<span class="l" id="L1086"><span class="n">1086</span>at every other prime.  Comparing local generating functions gives exactly</span>
<span class="l" id="L1087"><span class="n">1087</span>\[</span>
<span class="l" id="L1088"><span class="n">1088</span> v_{q,h}=\omega_q*c_{q,h}.</span>
<span class="l" id="L1089"><span class="n">1089</span>\]</span>
<span class="l" id="L1090"><span class="n">1090</span>Moreover</span>
<span class="l" id="L1091"><span class="n">1091</span>\[</span>
<span class="l" id="L1092"><span class="n">1092</span>1+\sum_{\nu\ge1}|c_{q,h}(p^\nu)|</span>
<span class="l" id="L1093"><span class="n">1093</span>=1+\frac1{(p-2)(p-4)},</span>
<span class="l" id="L1094"><span class="n">1094</span>\qquad</span>
<span class="l" id="L1095"><span class="n">1095</span>\sum_{\nu\ge1}|c_{q,h}(p^\nu)|\,\nu\log p</span>
<span class="l" id="L1096"><span class="n">1096</span>\ll\frac{\log p}{p^2},</span>
<span class="l" id="L1097"><span class="n">1097</span>\]</span>
<span class="l" id="L1098"><span class="n">1098</span>so the convergent product over active high primes gives</span>
<span class="l" id="L1099"><span class="n">1099</span>\begin{equation}\label{eq:fixed-q-convolution-transfer}</span>
<span class="l" id="L1100"><span class="n">1100</span> \sup_h\sum_s|c_{q,h}(s)|(1+\log s)&lt;\infty.</span>
<span class="l" id="L1101"><span class="n">1101</span>\end{equation}</span>
<span class="l" id="L1102"><span class="n">1102</span>Deleting the fixed primes dividing $q$ can only decrease these local norms.</span>
<span class="l" id="L1103"><span class="n">1103</span></span>
<span class="l" id="L1104"><span class="n">1104</span>For $v=v_{q,h}$ expand the convolution inside the partial sum.  Since every</span>
<span class="l" id="L1105"><span class="n">1105</span>$s$ in the support of $c_{q,h}$ is a unit modulo $\Lambda_q$,</span>
<span class="l" id="L1106"><span class="n">1106</span>\begin{align*}</span>
<span class="l" id="L1107"><span class="n">1107</span>&amp;\sum_{e\le X}\chi(e)v_{q,h}(e)</span>
<span class="l" id="L1108"><span class="n">1108</span>\psi\!\left(\frac{X_0}{\Lambda_q e}+\frac r{\Lambda_q}\right)\\</span>
<span class="l" id="L1109"><span class="n">1109</span>&amp;\qquad=\sum_{s\le X}c_{q,h}(s)\chi(s)</span>
<span class="l" id="L1110"><span class="n">1110</span>\sum_{m\le X/s}\chi(m)\omega_q(m)</span>
<span class="l" id="L1111"><span class="n">1111</span>\psi\!\left(\frac{X_0/s}{\Lambda_q m}+\frac r{\Lambda_q}\right).</span>
<span class="l" id="L1112"><span class="n">1112</span>\end{align*}</span>
<span class="l" id="L1113"><span class="n">1113</span>Apply \cref{lem:fixed-conductor-petermann} to the inner base sum with the real</span>
<span class="l" id="L1114"><span class="n">1114</span>reciprocal numerator $T=X_0/s$.  When $T&lt;2\Lambda_q$ the inner sum is</span>
<span class="l" id="L1115"><span class="n">1115</span>$O_q(1)$; otherwise \eqref{eq:petermann-real-numerator} applies.  Its right-hand</span>
<span class="l" id="L1116"><span class="n">1116</span>side is bounded by the same increasing majorant with $T$ replaced by $X_0$.</span>
<span class="l" id="L1117"><span class="n">1117</span>The uniform $\ell^1$ bound \eqref{eq:fixed-q-convolution-transfer} therefore</span>
<span class="l" id="L1118"><span class="n">1118</span>transfers that majorant without any integrality assumption on $X_0/s$.</span>
<span class="l" id="L1119"><span class="n">1119</span>For $\omega_q$ and $\eta_q$ take $s=1$.  Hence, for every coefficient</span>
<span class="l" id="L1120"><span class="n">1120</span>$v\in\{\omega_q,\eta_q,v_{q,h}\}$,</span>
<span class="l" id="L1121"><span class="n">1121</span>\begin{equation}\label{eq:petermann-q}</span>
<span class="l" id="L1122"><span class="n">1122</span>\begin{aligned}</span>
<span class="l" id="L1123"><span class="n">1123</span>&amp;\sum_{e\le X}\chi(e)v(e)</span>
<span class="l" id="L1124"><span class="n">1124</span>\psi\!\left(\frac{X_0}{\Lambda_q e}+\frac r{\Lambda_q}\right)\\</span>
<span class="l" id="L1125"><span class="n">1125</span>&amp;\qquad=O_q\!\left(</span>
<span class="l" id="L1126"><span class="n">1126</span>(\log X_0)^{3/4}</span>
<span class="l" id="L1127"><span class="n">1127</span>+(\log X_0)^{2/3}(\log\log X_0)^{4/3}</span>
<span class="l" id="L1128"><span class="n">1128</span>+\log\log X_0+1</span>
<span class="l" id="L1129"><span class="n">1129</span>\right)</span>
<span class="l" id="L1130"><span class="n">1130</span>=o_q(\log X_0).</span>
<span class="l" id="L1131"><span class="n">1131</span>\end{aligned}</span>
<span class="l" id="L1132"><span class="n">1132</span>\end{equation}</span>
<span class="l" id="L1133"><span class="n">1133</span>uniformly for $1\le X\le X_0$.</span>
<span class="l" id="L1134"><span class="n">1134</span></span>
<span class="l" id="L1135"><span class="n">1135</span>The principal periodic coefficient can also be checked locally.  For a residue</span>
<span class="l" id="L1136"><span class="n">1136</span>$r\pmod q$ define the free one-anchor low-period weight</span>
<span class="l" id="L1137"><span class="n">1137</span>\begin{equation}\label{eq:rho-free-q}</span>
<span class="l" id="L1138"><span class="n">1138</span>\rho_{q,r}(y)</span>
<span class="l" id="L1139"><span class="n">1139</span>=\1_{(r+y,q)=1}</span>
<span class="l" id="L1140"><span class="n">1140</span>\prod_{\substack{p\in\{2,3\}\\p\nmid q}}L_p(\{0,y\}).</span>
<span class="l" id="L1141"><span class="n">1141</span>\end{equation}</span>
<span class="l" id="L1142"><span class="n">1142</span></span>
<span class="l" id="L1143"><span class="n">1143</span>Put $h=2H$, $X_0=H-1$, and $\Lambda=\Lambda_q$.  The exact finite-period</span>
<span class="l" id="L1144"><span class="n">1144</span>floor decomposition needed for the one-dimensional bracket is as follows.</span>
<span class="l" id="L1145"><span class="n">1145</span>Define the two parity-reduced low weights</span>
<span class="l" id="L1146"><span class="n">1146</span>\begin{equation}\label{eq:one-d-reduced-weights-q}</span>
<span class="l" id="L1147"><span class="n">1147</span>w^{A}_{r,h}(z):=\lambda_{q,r,h}(2z),</span>
<span class="l" id="L1148"><span class="n">1148</span>\qquad</span>
<span class="l" id="L1149"><span class="n">1149</span>w^{F}_{r}(z):=\rho_{q,r}(2z),</span>
<span class="l" id="L1150"><span class="n">1150</span>\end{equation}</span>
<span class="l" id="L1151"><span class="n">1151</span>viewed as $\Lambda$-periodic functions.  For</span>
<span class="l" id="L1152"><span class="n">1152</span>$u\in(\mathbb Z/\Lambda\mathbb Z)^\times$ and $c\bmod\Lambda$, let</span>
<span class="l" id="L1153"><span class="n">1153</span>$\alpha_u(c)\in\{1,\ldots,\Lambda\}$ be the positive representative of</span>
<span class="l" id="L1154"><span class="n">1154</span>$c u^{-1}\pmod\Lambda$, with the zero class represented by $\Lambda$, and set</span>
<span class="l" id="L1155"><span class="n">1155</span>\begin{align}</span>
<span class="l" id="L1156"><span class="n">1156</span>W(w)&amp;:=\sum_{c\bmod\Lambda}w(c),\label{eq:one-d-W}\\</span>
<span class="l" id="L1157"><span class="n">1157</span>B_w(u)&amp;:=\sum_{c\bmod\Lambda}w(c)</span>
<span class="l" id="L1158"><span class="n">1158</span>\left(\frac12-\frac{\alpha_u(c)}\Lambda\right),\label{eq:one-d-B}\\</span>
<span class="l" id="L1159"><span class="n">1159</span>\Psi_w(u;t)&amp;:=\sum_{c\bmod\Lambda}w(c)</span>
<span class="l" id="L1160"><span class="n">1160</span>\psi\!\left(\frac t\Lambda+1-\frac{\alpha_u(c)}\Lambda\right).</span>
<span class="l" id="L1161"><span class="n">1161</span>\label{eq:one-d-Psi}</span>
<span class="l" id="L1162"><span class="n">1162</span>\end{align}</span>
<span class="l" id="L1163"><span class="n">1163</span>Then, for every $e$ on the support of the high coefficients and</span>
<span class="l" id="L1164"><span class="n">1164</span>$u_e\equiv e\pmod\Lambda$,</span>
<span class="l" id="L1165"><span class="n">1165</span>\begin{equation}\label{eq:exact-periodic-floor-count}</span>
<span class="l" id="L1166"><span class="n">1166</span>\boxed{</span>
<span class="l" id="L1167"><span class="n">1167</span>\sum_{1\le n\le X_0/e}w(en)</span>
<span class="l" id="L1168"><span class="n">1168</span>=\frac{W(w)}\Lambda\frac{X_0}{e}</span>
<span class="l" id="L1169"><span class="n">1169</span>+B_w(u_e)-\Psi_w(u_e;X_0/e).</span>
<span class="l" id="L1170"><span class="n">1170</span>}</span>
<span class="l" id="L1171"><span class="n">1171</span>\end{equation}</span>
<span class="l" id="L1172"><span class="n">1172</span>Indeed the contribution of the residue $c$ is exactly</span>
<span class="l" id="L1173"><span class="n">1173</span>\[</span>
<span class="l" id="L1174"><span class="n">1174</span>w(c)\left\lfloor</span>
<span class="l" id="L1175"><span class="n">1175</span>\frac{X_0}{\Lambda e}+1-\frac{\alpha_{u_e}(c)}\Lambda</span>
<span class="l" id="L1176"><span class="n">1176</span>\right\rfloor,</span>
<span class="l" id="L1177"><span class="n">1177</span>\]</span>
<span class="l" id="L1178"><span class="n">1178</span>and \eqref{eq:exact-periodic-floor-count} is just</span>
<span class="l" id="L1179"><span class="n">1179</span>$\lfloor x\rfloor=x-\psi(x)-\tfrac12$ summed over $c$.</span>
<span class="l" id="L1180"><span class="n">1180</span></span>
<span class="l" id="L1181"><span class="n">1181</span>Let</span>
<span class="l" id="L1182"><span class="n">1182</span>\[</span>
<span class="l" id="L1183"><span class="n">1183</span>N^F_{q,r}(h;e):=</span>
<span class="l" id="L1184"><span class="n">1184</span>\sum_{\substack{0&lt;y&lt;h\\e\mid y}}\rho_{q,r}(y),</span>
<span class="l" id="L1185"><span class="n">1185</span>\qquad</span>
<span class="l" id="L1186"><span class="n">1186</span>\Delta^F_{q,r}(h;e):=N^F_{q,r}(h;e)-\frac{N_q(r;h)}e,</span>
<span class="l" id="L1187"><span class="n">1187</span>\]</span>
<span class="l" id="L1188"><span class="n">1188</span>and define the residue-$r$ corner and the free high-prime baseline by</span>
<span class="l" id="L1189"><span class="n">1189</span>\[</span>
<span class="l" id="L1190"><span class="n">1190</span>C_q^{(r)}:=C_q^*(h)\Delta_{q,r}(h;1,1),</span>
<span class="l" id="L1191"><span class="n">1191</span>\qquad</span>
<span class="l" id="L1192"><span class="n">1192</span>C_{f,q}:=\prod_{\substack{p\ge5\\p\nmid q}}c_1(p).</span>
<span class="l" id="L1193"><span class="n">1193</span>\]</span>
<span class="l" id="L1194"><span class="n">1194</span>The exact high-prime Euler expansions give</span>
<span class="l" id="L1195"><span class="n">1195</span>\begin{align}</span>
<span class="l" id="L1196"><span class="n">1196</span>Q_{1,q}^{(r)}&amp;:=C_q^{(r)}+P_{1,q}^{(r)}</span>
<span class="l" id="L1197"><span class="n">1197</span>=C_q^*(h)\sum_{e\ge1}v_{q,h}(e)\Delta_{q,r}(h;e,1),</span>
<span class="l" id="L1198"><span class="n">1198</span>\label{eq:exact-Q1q}\\</span>
<span class="l" id="L1199"><span class="n">1199</span>S_{q,r}&amp;=C_{f,q}\sum_{e\ge1}\eta_q(e)\Delta^F_{q,r}(h;e).</span>
<span class="l" id="L1200"><span class="n">1200</span>\label{eq:exact-Sq}</span>
<span class="l" id="L1201"><span class="n">1201</span>\end{align}</span>
<span class="l" id="L1202"><span class="n">1202</span>For \eqref{eq:exact-Sq}, the subtraction is exact because primewise</span>
<span class="l" id="L1203"><span class="n">1203</span>\[</span>
<span class="l" id="L1204"><span class="n">1204</span>c_1(p)\left(1+\frac1{p(p-2)}\right)=1,</span>
<span class="l" id="L1205"><span class="n">1205</span>\]</span>
<span class="l" id="L1206"><span class="n">1206</span>so $C_{f,q}\sum_e\eta_q(e)/e=1$.  Also</span>
<span class="l" id="L1207"><span class="n">1207</span>$C_q^{(r)}=O_q(1)$ by the same complete-period mean used in</span>
<span class="l" id="L1208"><span class="n">1208</span>\cref{lem:fixed-q-row-reduction}.</span>
<span class="l" id="L1209"><span class="n">1209</span></span>
<span class="l" id="L1210"><span class="n">1210</span>Substituting \eqref{eq:exact-periodic-floor-count} into</span>
<span class="l" id="L1211"><span class="n">1211</span>\eqref{eq:exact-Q1q}--\eqref{eq:exact-Sq} gives the literal formulas</span>
<span class="l" id="L1212"><span class="n">1212</span>\begin{align}</span>
<span class="l" id="L1213"><span class="n">1213</span>Q_{1,q}^{(r)}</span>
<span class="l" id="L1214"><span class="n">1214</span>={}&amp;C_q^*(h)\sum_{e\le X_0}v_{q,h}(e)</span>
<span class="l" id="L1215"><span class="n">1215</span>\left\{</span>
<span class="l" id="L1216"><span class="n">1216</span>\frac{X_0W(w^A_{r,h})/\Lambda-N_q(r;h)}e</span>
<span class="l" id="L1217"><span class="n">1217</span>+B_{w^A_{r,h}}(u_e)-\Psi_{w^A_{r,h}}(u_e;X_0/e)</span>
<span class="l" id="L1218"><span class="n">1218</span>\right\}\notag\\</span>
<span class="l" id="L1219"><span class="n">1219</span>&amp;-C_q^*(h)N_q(r;h)\sum_{e&gt;X_0}\frac{v_{q,h}(e)}e,</span>
<span class="l" id="L1220"><span class="n">1220</span>\label{eq:exact-Q1-floor-sawtooth}\\</span>
<span class="l" id="L1221"><span class="n">1221</span>S_{q,r}</span>
<span class="l" id="L1222"><span class="n">1222</span>={}&amp;C_{f,q}\sum_{e\le X_0}\eta_q(e)</span>
<span class="l" id="L1223"><span class="n">1223</span>\left\{</span>
<span class="l" id="L1224"><span class="n">1224</span>\frac{X_0W(w^F_r)/\Lambda-N_q(r;h)}e</span>
<span class="l" id="L1225"><span class="n">1225</span>+B_{w^F_r}(u_e)-\Psi_{w^F_r}(u_e;X_0/e)</span>
<span class="l" id="L1226"><span class="n">1226</span>\right\}\notag\\</span>
<span class="l" id="L1227"><span class="n">1227</span>&amp;-C_{f,q}N_q(r;h)\sum_{e&gt;X_0}\frac{\eta_q(e)}e.</span>
<span class="l" id="L1228"><span class="n">1228</span>\label{eq:exact-S-floor-sawtooth}</span>
<span class="l" id="L1229"><span class="n">1229</span>\end{align}</span>
<span class="l" id="L1230"><span class="n">1230</span>These are identities, not asymptotic replacements.</span>
<span class="l" id="L1231"><span class="n">1231</span></span>
<span class="l" id="L1232"><span class="n">1232</span>For completeness, they decompose into exactly the analytic pieces bounded</span>
<span class="l" id="L1233"><span class="n">1233</span>above.  On the unit group,</span>
<span class="l" id="L1234"><span class="n">1234</span>\[</span>
<span class="l" id="L1235"><span class="n">1235</span>\1_{u_e=u}=\frac1{\ph(\Lambda)}</span>
<span class="l" id="L1236"><span class="n">1236</span>\sum_{\chi\bmod\Lambda}\overline{\chi(u)}\chi(e),</span>
<span class="l" id="L1237"><span class="n">1237</span>\]</span>
<span class="l" id="L1238"><span class="n">1238</span>so the $B$-terms are a finite family of fixed-character coefficient sums and</span>
<span class="l" id="L1239"><span class="n">1239</span>the $\Psi$-terms are a finite family of the sawtooth sums</span>
<span class="l" id="L1240"><span class="n">1240</span>\eqref{eq:petermann-periodic}.  The nonprincipal fixed-character coefficient</span>
<span class="l" id="L1241"><span class="n">1241</span>sums are $O_q(1)$; the principal one has</span>
<span class="l" id="L1242"><span class="n">1242</span>\[</span>
<span class="l" id="L1243"><span class="n">1243</span>\sum_{e\le X}\chi_0(e)v(e)=F_0(v)\log X+O_q(1)</span>
<span class="l" id="L1244"><span class="n">1244</span>\]</span>
<span class="l" id="L1245"><span class="n">1245</span>for $v=\eta_q$ and, uniformly in $h$, for $v=v_{q,h}$ by the exact convolution</span>
<span class="l" id="L1246"><span class="n">1246</span>\eqref{eq:fixed-q-convolution-transfer}.  The density terms in braces are</span>
<span class="l" id="L1247"><span class="n">1247</span>$O_q(1)/e$, since</span>
<span class="l" id="L1248"><span class="n">1248</span>\[</span>
<span class="l" id="L1249"><span class="n">1249</span>\frac{W(w^A_{r,h})}{\Lambda}</span>
<span class="l" id="L1250"><span class="n">1250</span>=\frac{W(w^F_r)}{\Lambda}=\frac{2\ph(q)}q,</span>
<span class="l" id="L1251"><span class="n">1251</span>\qquad</span>
<span class="l" id="L1252"><span class="n">1252</span>N_q(r;h)=\frac{2\ph(q)}qH+O_q(1),</span>
<span class="l" id="L1253"><span class="n">1253</span>\]</span>
<span class="l" id="L1254"><span class="n">1254</span>and the two tails in</span>
<span class="l" id="L1255"><span class="n">1255</span>\eqref{eq:exact-Q1-floor-sawtooth}--\eqref{eq:exact-S-floor-sawtooth} are</span>
<span class="l" id="L1256"><span class="n">1256</span>$O_q(1)$ by the corresponding logarithmic mean and partial summation.</span>
<span class="l" id="L1257"><span class="n">1257</span></span>
<span class="l" id="L1258"><span class="n">1258</span>It remains only to identify the complete principal boundary coefficient.  Put</span>
<span class="l" id="L1259"><span class="n">1259</span>\[</span>
<span class="l" id="L1260"><span class="n">1260</span>\overline B(w):=\frac1{\ph(\Lambda)}</span>
<span class="l" id="L1261"><span class="n">1261</span>\sum_{u\in(\mathbb Z/\Lambda\mathbb Z)^\times}B_w(u).</span>
<span class="l" id="L1262"><span class="n">1262</span>\]</span>
<span class="l" id="L1263"><span class="n">1263</span>For $c\ne0\pmod\Lambda$, the orbit of $cu^{-1}$ is the unit group modulo</span>
<span class="l" id="L1264"><span class="n">1264</span>$\Lambda/(c,\Lambda)$; pairing $x$ with its negative shows that the average of</span>
<span class="l" id="L1265"><span class="n">1265</span>$\alpha_u(c)/\Lambda$ is $1/2$.  For $c=0$ it is identically $1$.  Therefore</span>
<span class="l" id="L1266"><span class="n">1266</span>\begin{equation}\label{eq:Bbar-zero-residue}</span>
<span class="l" id="L1267"><span class="n">1267</span>\boxed{\overline B(w)=-\frac12w(0).}</span>
<span class="l" id="L1268"><span class="n">1268</span>\end{equation}</span>
<span class="l" id="L1269"><span class="n">1269</span>At the zero low residue the anchored and free weights agree exactly:</span>
<span class="l" id="L1270"><span class="n">1270</span>for every active $p\in\{2,3\}$,</span>
<span class="l" id="L1271"><span class="n">1271</span>\[</span>
<span class="l" id="L1272"><span class="n">1272</span>\frac{L_p(\{0,0,h\})}{L_p(\{0,h\})}</span>
<span class="l" id="L1273"><span class="n">1273</span>=(1-1/p)^{-1}=L_p(\{0,0\}),</span>
<span class="l" id="L1274"><span class="n">1274</span>\]</span>
<span class="l" id="L1275"><span class="n">1275</span>and the $q$-admissibility indicator is $1$.  Hence</span>
<span class="l" id="L1276"><span class="n">1276</span>\begin{equation}\label{eq:Bbar-anchored-free}</span>
<span class="l" id="L1277"><span class="n">1277</span>\overline B(w^A_{r,h})=\overline B(w^F_r).</span>
<span class="l" id="L1278"><span class="n">1278</span>\end{equation}</span>
<span class="l" id="L1279"><span class="n">1279</span>CRT and the low-prime local tables give the exact complete-period means</span>
<span class="l" id="L1280"><span class="n">1280</span>\begin{equation}\label{eq:principal-low-mean}</span>
<span class="l" id="L1281"><span class="n">1281</span>\frac1{\Lambda_q}\sum_{c\bmod\Lambda_q}\lambda_{q,r,h}(c)</span>
<span class="l" id="L1282"><span class="n">1282</span>=\frac1{\Lambda_q}\sum_{c\bmod\Lambda_q}\rho_{q,r}(c)</span>
<span class="l" id="L1283"><span class="n">1283</span>=\frac{\ph(q)}q.</span>
<span class="l" id="L1284"><span class="n">1284</span>\end{equation}</span>
<span class="l" id="L1285"><span class="n">1285</span>Indeed every prime $p\mid q$ contributes the common factor $1-1/p$; at an</span>
<span class="l" id="L1286"><span class="n">1286</span>active prime $2$ the two local weights have values $(2,0)$ and mean one; at</span>
<span class="l" id="L1287"><span class="n">1287</span>an active prime $3$ the free values are $(3/2,3/4,3/4)$, while the anchored</span>
<span class="l" id="L1288"><span class="n">1288</span>values are the same when $3\mid h$ and, when $3\nmid h$, are a permutation</span>
<span class="l" id="L1289"><span class="n">1289</span>of $(3/2,3/2,0)$.  Each active low-prime mean is therefore one.</span>
<span class="l" id="L1290"><span class="n">1290</span></span>
<span class="l" id="L1291"><span class="n">1291</span>For a squarefree coefficient $u$ occurring here, write its logarithmic Euler</span>
<span class="l" id="L1292"><span class="n">1292</span>residue as</span>
<span class="l" id="L1293"><span class="n">1293</span>\[</span>
<span class="l" id="L1294"><span class="n">1294</span>F_0(u):=\prod_p\left(1-\frac1p\right)(1+u(p)),</span>
<span class="l" id="L1295"><span class="n">1295</span>\]</span>
<span class="l" id="L1296"><span class="n">1296</span>and recall the free high-prime baseline $C_{f,q}$ defined above.</span>
<span class="l" id="L1297"><span class="n">1297</span>The active high-prime identities</span>
<span class="l" id="L1298"><span class="n">1298</span>\[</span>
<span class="l" id="L1299"><span class="n">1299</span> c_0(p)\left(1-\frac1p\right)\left(1+\frac1{p-3}\right)=1,</span>
<span class="l" id="L1300"><span class="n">1300</span> \qquad</span>
<span class="l" id="L1301"><span class="n">1301</span> c_1(p)\left(1-\frac1p\right)\left(1+\frac1{p-2}\right)=1</span>
<span class="l" id="L1302"><span class="n">1302</span>\]</span>
<span class="l" id="L1303"><span class="n">1303</span>then give the literal coefficient equality</span>
<span class="l" id="L1304"><span class="n">1304</span>\begin{equation}\label{eq:principal-high-coefficient}</span>
<span class="l" id="L1305"><span class="n">1305</span>C_q^*(h)F_0(v_{q,h})</span>
<span class="l" id="L1306"><span class="n">1306</span>=C_{f,q}F_0(\eta_q)</span>
<span class="l" id="L1307"><span class="n">1307</span>=\prod_{p\mid6q}\left(1-\frac1p\right)</span>
<span class="l" id="L1308"><span class="n">1308</span>=\frac{\ph(\Lambda_q)}{\Lambda_q}.</span>
<span class="l" id="L1309"><span class="n">1309</span>\end{equation}</span>
<span class="l" id="L1310"><span class="n">1310</span>The exact floor/sawtooth decompositions</span>
<span class="l" id="L1311"><span class="n">1311</span>\eqref{eq:exact-Q1-floor-sawtooth}--\eqref{eq:exact-S-floor-sawtooth} now leave</span>
<span class="l" id="L1312"><span class="n">1312</span>no unassigned one-dimensional term.  Their finite sawtooth character families, including the principal character,</span>
<span class="l" id="L1313"><span class="n">1313</span>are $o_q(\log h)$ by \eqref{eq:petermann-q}; the nonprincipal</span>
<span class="l" id="L1314"><span class="n">1314</span>constant-character pieces and all density, endpoint, and tail terms are</span>
<span class="l" id="L1315"><span class="n">1315</span>$O_q(1)$.  The only logarithmic principal terms are</span>
<span class="l" id="L1316"><span class="n">1316</span>\[</span>
<span class="l" id="L1317"><span class="n">1317</span>C_q^*(h)F_0(v_{q,h})\overline B(w^A_{r,h})\log X_0</span>
<span class="l" id="L1318"><span class="n">1318</span>\quad\text{and}\quad</span>
<span class="l" id="L1319"><span class="n">1319</span>C_{f,q}F_0(\eta_q)\overline B(w^F_r)\log X_0.</span>
<span class="l" id="L1320"><span class="n">1320</span>\]</span>
<span class="l" id="L1321"><span class="n">1321</span>They are equal by \eqref{eq:principal-high-coefficient} and</span>
<span class="l" id="L1322"><span class="n">1322</span>\eqref{eq:Bbar-anchored-free}.  Thus</span>
<span class="l" id="L1323"><span class="n">1323</span>$Q_{1,q}^{(r)}-S_{q,r}=o_q(\log h)$; since</span>
<span class="l" id="L1324"><span class="n">1324</span>$C_q^{(r)}=O_q(1)$, the residue-sensitive notation</span>
<span class="l" id="L1325"><span class="n">1325</span>\eqref{eq:P1-residue-q} gives</span>
<span class="l" id="L1326"><span class="n">1326</span>\begin{equation}\label{eq:pure-one-uniform-q}</span>
<span class="l" id="L1327"><span class="n">1327</span>P_{1,q}^{(r)}-S_{q,r}=o_q(\log h)</span>
<span class="l" id="L1328"><span class="n">1328</span>\end{equation}</span>
<span class="l" id="L1329"><span class="n">1329</span>for every compatible unit residue $r\pmod q$ occurring in the reflected pair.</span>
<span class="l" id="L1330"><span class="n">1330</span>Taking $r=a$ proves the first bracket.  Taking $r=-b$ and then applying the</span>
<span class="l" id="L1331"><span class="n">1331</span>exact family reflection \eqref{eq:reflection-assembly-q} and support bound</span>
<span class="l" id="L1332"><span class="n">1332</span>\eqref{eq:support-bound-q} proves the second.  Explicitly,</span>
<span class="l" id="L1333"><span class="n">1333</span>\begin{equation}\label{eq:one-dimensional-q}</span>
<span class="l" id="L1334"><span class="n">1334</span>P_{1,q}^{(a)}-S_{q,a}=o_q(\log h),</span>
<span class="l" id="L1335"><span class="n">1335</span>\qquad</span>
<span class="l" id="L1336"><span class="n">1336</span>P_{2,q}+D_{1,q}-S_{q,-b}</span>
<span class="l" id="L1337"><span class="n">1337</span>=\bigl[P_{1,q}^{(-b)}-S_{q,-b}\bigr]</span>
<span class="l" id="L1338"><span class="n">1338</span>-P_{1,q,h}^{\mathrm{supp}}(-b)</span>
<span class="l" id="L1339"><span class="n">1339</span>=o_q(\log h).</span>
<span class="l" id="L1340"><span class="n">1340</span>\end{equation}</span>
<span class="l" id="L1341"><span class="n">1341</span></span>
<span class="l" id="L1342"><span class="n">1342</span>\section{Proof of the fixed-modulus theorem}\label{sec:proof-main}</span>
<span class="l" id="L1343"><span class="n">1343</span></span>
<span class="l" id="L1344"><span class="n">1344</span>Only assembly remains.  The exact divisor expansion \eqref{eq:Tq-expansion} has four geometrically different pieces: the corner, two one-dimensional boundary families, the ordinary two-dimensional rows, and the rows decorated by height primes.  Let $C_q$ be the corner $(e_1,e_2)=(1,1)$; let $P_{1,q}=P_{1,q}^{(a)}$ and $P_{2,q}$ denote respectively the rows $e_2=1&lt;e_1$ and $e_1=1&lt;e_2$; let $G_q$ denote the genuinely two-dimensional rows $e_1,e_2&gt;1$ with $(e_1,h)=1$; and let $D_q$ denote those with $(e_1,h)&gt;1$.  Write $e_1=dr$, where $d=(e_1,h)$.  Split</span>
<span class="l" id="L1345"><span class="n">1345</span>\[</span>
<span class="l" id="L1346"><span class="n">1346</span>D_q=D_{1,q}+D_{2D,q}</span>
<span class="l" id="L1347"><span class="n">1347</span>\]</span>
<span class="l" id="L1348"><span class="n">1348</span>according as $r=1$ or $r&gt;1$.  Only $D_{1,q}$ is one-dimensional; $D_{2D,q}$ is a decorated two-dimensional family handled by \eqref{eq:decorated-map}.</span>
<span class="l" id="L1349"><span class="n">1349</span></span>
<span class="l" id="L1350"><span class="n">1350</span>The exact family assembly is</span>
<span class="l" id="L1351"><span class="n">1351</span>\begin{equation}\label{eq:final-assembly-q}</span>
<span class="l" id="L1352"><span class="n">1352</span>\boxed{</span>
<span class="l" id="L1353"><span class="n">1353</span>A_{1,q}-B_{1,q}</span>
<span class="l" id="L1354"><span class="n">1354</span>=[P_{1,q}-S_{q,a}]</span>
<span class="l" id="L1355"><span class="n">1355</span>+[P_{2,q}+D_{1,q}-S_{q,-b}]</span>
<span class="l" id="L1356"><span class="n">1356</span>+C_q+[G_q+D_{2D,q}].</span>
<span class="l" id="L1357"><span class="n">1357</span>}</span>
<span class="l" id="L1358"><span class="n">1358</span>\end{equation}</span>
<span class="l" id="L1359"><span class="n">1359</span>The corner and the exact periodic-density corrections from</span>
<span class="l" id="L1360"><span class="n">1360</span>\eqref{eq:fixed-q-row-reduction} contribute</span>
<span class="l" id="L1361"><span class="n">1361</span>\[</span>
<span class="l" id="L1362"><span class="n">1362</span>O_q(1)+O_q(\log\log h).</span>
<span class="l" id="L1363"><span class="n">1363</span>\]</span>
<span class="l" id="L1364"><span class="n">1364</span>The two one-dimensional brackets are $o_q(\log h)$ by \eqref{eq:one-dimensional-q}.  The centered two-dimensional small box is $o_q(\log h)$ by \eqref{eq:fixed-q-smallbox}--\eqref{eq:exponent}; its hyperbolic fill and the full exterior are $o_q(\log h)$ by \eqref{eq:exterior-q}; and the decorated rows remain $o_q(\log h)$ by \eqref{eq:decorated-map}.  Inserting these estimates into \eqref{eq:final-assembly-q} proves \cref{thm:fixed-q}.</span>
<span class="l" id="L1365"><span class="n">1365</span></span>
<span class="l" id="L1366"><span class="n">1366</span>\begin{remark}[Asymptotic character]\label{rem:asymptotic}</span>
<span class="l" id="L1367"><span class="n">1367</span>The strict Boolean margin is inherited from the $q=1$ proof and is numerically small.  The reciprocal argument also uses a superlogarithmic but very slowly emerging saving.  The theorem is therefore asymptotic; no computationally useful finite-height threshold is asserted.</span>
<span class="l" id="L1368"><span class="n">1368</span>\end{remark}</span>
<span class="l" id="L1369"><span class="n">1369</span></span>
<span class="l" id="L1370"><span class="n">1370</span>\section{Finite-cutoff arithmetic diagonal contraction}\label{sec:collision}</span>
<span class="l" id="L1371"><span class="n">1371</span></span>
<span class="l" id="L1372"><span class="n">1372</span>The level-one transfer theorem is already proved.  This section records a separate finite-cutoff identity for the active high-prime coefficient system.  Boolean M\&quot;obius inversion is an incidence-algebra operation; by contrast, averaging a repeated free coordinate after connectedization produces an arithmetic operator on the modulus coefficients.</span>
<span class="l" id="L1373"><span class="n">1373</span></span>
<span class="l" id="L1374"><span class="n">1374</span>Fix a finite active high-prime set $\mathcal P$ and an admissible anchor pattern</span>
<span class="l" id="L1375"><span class="n">1375</span>$A&#x27;\subseteq E$.  Free coordinates remain labeled even when two coordinate</span>
<span class="l" id="L1376"><span class="n">1376</span>values coincide.  For $p\in\mathcal P$ define</span>
<span class="l" id="L1377"><span class="n">1377</span>\begin{equation}\label{eq:collision-local}</span>
<span class="l" id="L1378"><span class="n">1378</span>\phi_{p,z}(x)=\frac{\1_{x\not\equiv z\ (p)}}{1-1/p},</span>
<span class="l" id="L1379"><span class="n">1379</span>\qquad</span>
<span class="l" id="L1380"><span class="n">1380</span>L_p(\mathbf x)=\frac1p\sum_{z\bmod p}\prod_j\phi_{p,z}(x_j),</span>
<span class="l" id="L1381"><span class="n">1381</span>\end{equation}</span>
<span class="l" id="L1382"><span class="n">1382</span>and</span>
<span class="l" id="L1383"><span class="n">1383</span>\begin{equation}\label{eq:collision-gw}</span>
<span class="l" id="L1384"><span class="n">1384</span>g_p^{A&#x27;}(\mathbf x)=\frac{L_p(A&#x27;\cup\mathbf x)}{L_p(A&#x27;)},</span>
<span class="l" id="L1385"><span class="n">1385</span>\qquad</span>
<span class="l" id="L1386"><span class="n">1386</span>w_p^{A&#x27;}(\mathbf x)=g_p^{A&#x27;}(\mathbf x)-1.</span>
<span class="l" id="L1387"><span class="n">1387</span>\end{equation}</span>
<span class="l" id="L1388"><span class="n">1388</span>For squarefree $m\mid\prod_{p\in\mathcal P}p$ set</span>
<span class="l" id="L1389"><span class="n">1389</span>\begin{equation}\label{eq:collision-u}</span>
<span class="l" id="L1390"><span class="n">1390</span>u_1^{A&#x27;}=1,</span>
<span class="l" id="L1391"><span class="n">1391</span>\qquad</span>
<span class="l" id="L1392"><span class="n">1392</span>u_m^{A&#x27;}(\mathbf x)=\prod_{p\mid m}w_p^{A&#x27;}(\mathbf x).</span>
<span class="l" id="L1393"><span class="n">1393</span>\end{equation}</span>
<span class="l" id="L1394"><span class="n">1394</span>Then the truncated relative Euler product is reconstructed exactly by</span>
<span class="l" id="L1395"><span class="n">1395</span>\begin{equation}\label{eq:collision-euler-reconstruction}</span>
<span class="l" id="L1396"><span class="n">1396</span>\sum_{m\mid\prod_{p\in\mathcal P}p}u_m^{A&#x27;}(\mathbf x)</span>
<span class="l" id="L1397"><span class="n">1397</span>=\prod_{p\in\mathcal P}g_p^{A&#x27;}(\mathbf x).</span>
<span class="l" id="L1398"><span class="n">1398</span>\end{equation}</span>
<span class="l" id="L1399"><span class="n">1399</span>For a labeled set $J$ define</span>
<span class="l" id="L1400"><span class="n">1400</span>\begin{equation}\label{eq:collision-mob}</span>
<span class="l" id="L1401"><span class="n">1401</span>\operatorname{Mob}_JF</span>
<span class="l" id="L1402"><span class="n">1402</span>=\sum_{V\subseteq J}(-1)^{|J|-|V|}F(\mathbf x_V).</span>
<span class="l" id="L1403"><span class="n">1403</span>\end{equation}</span>
<span class="l" id="L1404"><span class="n">1404</span>For a doubled block $B=\{1,2\}$ define diagonal contraction at modulus $m$ by</span>
<span class="l" id="L1405"><span class="n">1405</span>\begin{equation}\label{eq:collision-D}</span>
<span class="l" id="L1406"><span class="n">1406</span>(\mathsf D_BF)_m(\mathbf y)</span>
<span class="l" id="L1407"><span class="n">1407</span>=\frac1m\sum_{x\bmod m}F_m(x,x,\mathbf y).</span>
<span class="l" id="L1408"><span class="n">1408</span>\end{equation}</span>
<span class="l" id="L1409"><span class="n">1409</span>Finally put</span>
<span class="l" id="L1410"><span class="n">1410</span>\begin{equation}\label{eq:L-operator}</span>
<span class="l" id="L1411"><span class="n">1411</span>(\Kcal f)_m=\frac1{\ph(m)}\sum_{d\mid m}d f_d,</span>
<span class="l" id="L1412"><span class="n">1412</span>\qquad</span>
<span class="l" id="L1413"><span class="n">1413</span>\Lcal=\Kcal-I.</span>
<span class="l" id="L1414"><span class="n">1414</span>\end{equation}</span>
<span class="l" id="L1415"><span class="n">1415</span></span>
<span class="l" id="L1416"><span class="n">1416</span>\begin{proposition}[Finite-cutoff collision intertwiner]\label{prop:collision-intertwiner}</span>
<span class="l" id="L1417"><span class="n">1417</span>For every squarefree finite-cutoff modulus $m$,</span>
<span class="l" id="L1418"><span class="n">1418</span>\begin{equation}\label{eq:collision-intertwiner}</span>
<span class="l" id="L1419"><span class="n">1419</span>\boxed{</span>
<span class="l" id="L1420"><span class="n">1420</span>\mathsf D_B\operatorname{Mob}_{B\cup U}u_m^{A&#x27;}</span>
<span class="l" id="L1421"><span class="n">1421</span>=\operatorname{Mob}_U(\Lcal u^{A&#x27;})_m.</span>
<span class="l" id="L1422"><span class="n">1422</span>}</span>
<span class="l" id="L1423"><span class="n">1423</span>\end{equation}</span>
<span class="l" id="L1424"><span class="n">1424</span>\end{proposition}</span>
<span class="l" id="L1425"><span class="n">1425</span></span>
<span class="l" id="L1426"><span class="n">1426</span>\begin{proof}</span>
<span class="l" id="L1427"><span class="n">1427</span>Write</span>
<span class="l" id="L1428"><span class="n">1428</span>\[</span>
<span class="l" id="L1429"><span class="n">1429</span>G_d^{A&#x27;}(\mathbf x)=\prod_{p\mid d}g_p^{A&#x27;}(\mathbf x).</span>
<span class="l" id="L1430"><span class="n">1430</span>\]</span>
<span class="l" id="L1431"><span class="n">1431</span>Then inclusion--exclusion in the active prime set gives</span>
<span class="l" id="L1432"><span class="n">1432</span>\begin{equation}\label{eq:collision-uG}</span>
<span class="l" id="L1433"><span class="n">1433</span>u_m^{A&#x27;}</span>
<span class="l" id="L1434"><span class="n">1434</span>=\sum_{d\mid m}(-1)^{\Om(m/d)}G_d^{A&#x27;}.</span>
<span class="l" id="L1435"><span class="n">1435</span>\end{equation}</span>
<span class="l" id="L1436"><span class="n">1436</span>Under the anchor-conditioned product measure, $G_d^{A&#x27;}$ is an expectation</span>
<span class="l" id="L1437"><span class="n">1437</span>of products of</span>
<span class="l" id="L1438"><span class="n">1438</span>\[</span>
<span class="l" id="L1439"><span class="n">1439</span>\Phi_{d,\mathbf z}(x)=\prod_{p\mid d}\phi_{p,z_p}(x).</span>
<span class="l" id="L1440"><span class="n">1440</span>\]</span>
<span class="l" id="L1441"><span class="n">1441</span>Boolean connectedization therefore produces a product of</span>
<span class="l" id="L1442"><span class="n">1442</span>$\Phi_{d,\mathbf z}-1$.  CRT gives the exact diagonal second moment</span>
<span class="l" id="L1443"><span class="n">1443</span>\begin{equation}\label{eq:collision-second-moment}</span>
<span class="l" id="L1444"><span class="n">1444</span>\frac1d\sum_{x\bmod d}(\Phi_{d,\mathbf z}(x)-1)^2</span>
<span class="l" id="L1445"><span class="n">1445</span>=\frac d{\ph(d)}-1,</span>
<span class="l" id="L1446"><span class="n">1446</span>\end{equation}</span>
<span class="l" id="L1447"><span class="n">1447</span>independently of the anchor-conditioned $\mathbf z$.  Hence</span>
<span class="l" id="L1448"><span class="n">1448</span>\begin{equation}\label{eq:collision-G}</span>
<span class="l" id="L1449"><span class="n">1449</span>\mathsf D_B\operatorname{Mob}_{B\cup U}G_d^{A&#x27;}</span>
<span class="l" id="L1450"><span class="n">1450</span>=\left(\frac d{\ph(d)}-1\right)\operatorname{Mob}_UG_d^{A&#x27;}.</span>
<span class="l" id="L1451"><span class="n">1451</span>\end{equation}</span>
<span class="l" id="L1452"><span class="n">1452</span>Substituting \eqref{eq:collision-uG} into \eqref{eq:L-operator}, the coefficient</span>
<span class="l" id="L1453"><span class="n">1453</span>of $G_d^{A&#x27;}$ in $(\Lcal u^{A&#x27;})_m$ is exactly</span>
<span class="l" id="L1454"><span class="n">1454</span>\[</span>
<span class="l" id="L1455"><span class="n">1455</span>(-1)^{\Om(m/d)}\left(\frac d{\ph(d)}-1\right).</span>
<span class="l" id="L1456"><span class="n">1456</span>\]</span>
<span class="l" id="L1457"><span class="n">1457</span>Comparison with \eqref{eq:collision-G} proves</span>
<span class="l" id="L1458"><span class="n">1458</span>\eqref{eq:collision-intertwiner}.</span>
<span class="l" id="L1459"><span class="n">1459</span>\end{proof}</span>
<span class="l" id="L1460"><span class="n">1460</span></span>
<span class="l" id="L1461"><span class="n">1461</span>Define the anchor-combined rank-$r$ transfer kernel by</span>
<span class="l" id="L1462"><span class="n">1462</span>\begin{equation}\label{eq:collision-Dr}</span>
<span class="l" id="L1463"><span class="n">1463</span>D_m^{(r)}(\mathbf x)</span>
<span class="l" id="L1464"><span class="n">1464</span>=\operatorname{Mob}_{[r]}</span>
<span class="l" id="L1465"><span class="n">1465</span>\sum_{A&#x27;\subseteq E}(-1)^{|E|-|A&#x27;|}u_m^{A&#x27;}(\mathbf x).</span>
<span class="l" id="L1466"><span class="n">1466</span>\end{equation}</span>
<span class="l" id="L1467"><span class="n">1467</span>Because $\Lcal$ is independent of the anchor pattern, it commutes with this</span>
<span class="l" id="L1468"><span class="n">1468</span>anchor combination and gives</span>
<span class="l" id="L1469"><span class="n">1469</span>\begin{equation}\label{eq:collision-one}</span>
<span class="l" id="L1470"><span class="n">1470</span>\frac1m\sum_{x\bmod m}D_m^{(r+2)}(x,x,\mathbf y)</span>
<span class="l" id="L1471"><span class="n">1471</span>=(\Lcal D^{(r)})_m(\mathbf y).</span>
<span class="l" id="L1472"><span class="n">1472</span>\end{equation}</span>
<span class="l" id="L1473"><span class="n">1473</span>Iterating disjoint doubled blocks gives</span>
<span class="l" id="L1474"><span class="n">1474</span>\begin{equation}\label{eq:collision-iterate}</span>
<span class="l" id="L1475"><span class="n">1475</span>\frac1{m^j}\sum_{x_1,\ldots,x_j\bmod m}</span>
<span class="l" id="L1476"><span class="n">1476</span>D_m^{(2j+1)}(x_1,x_1,\ldots,x_j,x_j,y)</span>
<span class="l" id="L1477"><span class="n">1477</span>=(\Lcal^jD^{(1)})_m(y).</span>
<span class="l" id="L1478"><span class="n">1478</span>\end{equation}</span>
<span class="l" id="L1479"><span class="n">1479</span></span>
<span class="l" id="L1480"><span class="n">1480</span>Equations \eqref{eq:collision-intertwiner}--\eqref{eq:collision-iterate} are</span>
<span class="l" id="L1481"><span class="n">1481</span>finite-cutoff identities for the active high-prime system alone.  They are not</span>
<span class="l" id="L1482"><span class="n">1482</span>used in the proof of \cref{thm:fixed-q}, and no asymptotic or rank-uniform</span>
<span class="l" id="L1483"><span class="n">1483</span>statement is attached to them here.</span>
<span class="l" id="L1484"><span class="n">1484</span></span>
<span class="l" id="L1485"><span class="n">1485</span>\section{Relation to the completed fixed-gap density}\label{sec:los}</span>
<span class="l" id="L1486"><span class="n">1486</span></span>
<span class="l" id="L1487"><span class="n">1487</span>The contraction identity \eqref{eq:F-Rcirc} is the exact bridge to the completed endpoint coordinates used in the fixed-gap consecutive-prime residue formalism.  At level one,</span>
<span class="l" id="L1488"><span class="n">1488</span>\[</span>
<span class="l" id="L1489"><span class="n">1489</span>C_{1,q}^{\mathrm{comp}}=\SSer_q(E)A_{1,q},</span>
<span class="l" id="L1490"><span class="n">1490</span>\]</span>
<span class="l" id="L1491"><span class="n">1491</span>while the transfer theorem controls the connected defect between this anchored coordinate and the sum of the two one-anchor connected boundaries.  The ambient modulus density is carried separately by \eqref{eq:deltaq}; because $\delta_q\ph(q)/q=1/\log x$, there is no hidden $q/\ph(q)$ renormalization inside the connected transfer kernel.  Nothing in this note supplies a rank-uniform tail estimate or rederives a Poisson survival factor.</span>
<span class="l" id="L1492"><span class="n">1492</span></span>
<span class="l" id="L1493"><span class="n">1493</span>\section{Source and priority note}\label{sec:priority}</span>
<span class="l" id="L1494"><span class="n">1494</span></span>
<span class="l" id="L1495"><span class="n">1495</span>The mathematical assertions of \cref{prop:endpoint-contraction,thm:fixed-q} are established by the proofs above.  Priority standing is separate; proof standing is not used as evidence of novelty or priority.</span>
<span class="l" id="L1496"><span class="n">1496</span></span>
<span class="l" id="L1497"><span class="n">1497</span>We prove the exact contraction dictionary and the fixed-$q$ pointwise connected transfer theorem stated above. A 2026 preprint of Hamed Ebadi, referenced by a public Zenodo verification record \cite{Ebadi2026}, studies the same fixed-two-anchor raw coefficient hierarchy at $q=1$. The manuscript proof was not available to us for inspection at the time of this version, so we make no claim here about priority relative to proof-level lemmas in that preprint beyond the literal distinctions supported by its public theorem statement.</span>
<span class="l" id="L1498"><span class="n">1498</span></span>
<span class="l" id="L1499"><span class="n">1499</span>For the same reason, this note makes no novelty claim for fixed anchors, congruence restrictions, Boolean M\&quot;obius inversion, generic connected/cumulant language, or the conditional fixed-gap Poisson factor.  Those issues belong to the prior-art record rather than to the proof of \cref{thm:fixed-q}.</span>
<span class="l" id="L1500"><span class="n">1500</span></span>
<span class="l" id="L1501"><span class="n">1501</span>\begin{thebibliography}{9}</span>
<span class="l" id="L1502"><span class="n">1502</span></span>
<span class="l" id="L1503"><span class="n">1503</span>\bibitem{DorseyPointwise2026}</span>
<span class="l" id="L1504"><span class="n">1504</span>C.~Dorsey,</span>
<span class="l" id="L1505"><span class="n">1505</span>\emph{Boolean CRT Fibers and Pointwise Anchor Contraction},</span>
<span class="l" id="L1506"><span class="n">1506</span>preprint manuscript, September 2026.</span>
<span class="l" id="L1507"><span class="n">1507</span></span>
<span class="l" id="L1508"><span class="n">1508</span>\bibitem{ShiuThesis1980}</span>
<span class="l" id="L1509"><span class="n">1509</span>P.~M.-K.~Shiu,</span>
<span class="l" id="L1510"><span class="n">1510</span>\emph{Studies in Multiplicative Number Theory},</span>
<span class="l" id="L1511"><span class="n">1511</span>Ph.D. thesis, University of Nottingham, May 1980, Chapter~3.</span>
<span class="l" id="L1512"><span class="n">1512</span></span>
<span class="l" id="L1513"><span class="n">1513</span>\bibitem{Karatsuba1971}</span>
<span class="l" id="L1514"><span class="n">1514</span>A.~A.~Karatsuba,</span>
<span class="l" id="L1515"><span class="n">1515</span>\emph{Estimates for trigonometric sums by Vinogradov&#x27;s method, and some applications},</span>
<span class="l" id="L1516"><span class="n">1516</span>Proc. Steklov Inst. Math. 112 (1971), 251--265</span>
<span class="l" id="L1517"><span class="n">1517</span>(English translation, American Mathematical Society, 1973).</span>
<span class="l" id="L1518"><span class="n">1518</span></span>
<span class="l" id="L1519"><span class="n">1519</span>\bibitem{Petermann1998}</span>
<span class="l" id="L1520"><span class="n">1520</span>Y.-F.~S.~P\&#x27;etermann,</span>
<span class="l" id="L1521"><span class="n">1521</span>\emph{On an estimate of Walfisz and Saltykov for an error term related to the Euler function},</span>
<span class="l" id="L1522"><span class="n">1522</span>J. Th\&#x27;eor. Nombres Bordeaux 10 (1998), no.~1, 203--236.</span>
<span class="l" id="L1523"><span class="n">1523</span></span>
<span class="l" id="L1524"><span class="n">1524</span>\bibitem{GoldstonLedoan2011}</span>
<span class="l" id="L1525"><span class="n">1525</span>D.~A.~Goldston and A.~H.~Ledoan,</span>
<span class="l" id="L1526"><span class="n">1526</span>\emph{On the differences between consecutive prime numbers, I},</span>
<span class="l" id="L1527"><span class="n">1527</span>arXiv:1111.3380.</span>
<span class="l" id="L1528"><span class="n">1528</span></span>
<span class="l" id="L1529"><span class="n">1529</span>\bibitem{LOS}</span>
<span class="l" id="L1530"><span class="n">1530</span>R.~J.~Lemke Oliver and K.~Soundararajan,</span>
<span class="l" id="L1531"><span class="n">1531</span>\emph{Unexpected biases in the distribution of consecutive primes},</span>
<span class="l" id="L1532"><span class="n">1532</span>preprint manuscript.</span>
<span class="l" id="L1533"><span class="n">1533</span></span>
<span class="l" id="L1534"><span class="n">1534</span>\bibitem{Kuperberg2025}</span>
<span class="l" id="L1535"><span class="n">1535</span>V.~Kuperberg,</span>
<span class="l" id="L1536"><span class="n">1536</span>\emph{Sums of singular series along arithmetic progressions and with smooth weights},</span>
<span class="l" id="L1537"><span class="n">1537</span>arXiv:2301.06095v1 (2023).</span>
<span class="l" id="L1538"><span class="n">1538</span></span>
<span class="l" id="L1539"><span class="n">1539</span>\bibitem{Ebadi2026}</span>
<span class="l" id="L1540"><span class="n">1540</span>H.~Ebadi,</span>
<span class="l" id="L1541"><span class="n">1541</span>2026 preprint referenced by Zenodo verification record 21534407.</span>
<span class="l" id="L1542"><span class="n">1542</span>Manuscript proof unavailable to the author of this version at the time of writing.</span>
<span class="l" id="L1543"><span class="n">1543</span></span>
<span class="l" id="L1544"><span class="n">1544</span>\end{thebibliography}</span>
<span class="l" id="L1545"><span class="n">1545</span></span>
<span class="l" id="L1546"><span class="n">1546</span>\end{document}</span></pre></div></body></html>